% Options for packages loaded elsewhere
\PassOptionsToPackage{unicode}{hyperref}
\PassOptionsToPackage{hyphens}{url}
\documentclass[
  11pt,
]{article}
\usepackage{xcolor}
\usepackage[left=3cm,right=3cm,top=2.5cm,bottom=2.5cm]{geometry}
\usepackage{amsmath,amssymb}
\setcounter{secnumdepth}{5}
\usepackage{iftex}
\ifPDFTeX
  \usepackage[T1]{fontenc}
  \usepackage[utf8]{inputenc}
  \usepackage{textcomp} % provide euro and other symbols
\else % if luatex or xetex
  \usepackage{unicode-math} % this also loads fontspec
  \defaultfontfeatures{Scale=MatchLowercase}
  \defaultfontfeatures[\rmfamily]{Ligatures=TeX,Scale=1}
\fi
\usepackage{lmodern}
\ifPDFTeX\else
  % xetex/luatex font selection
\fi
% Use upquote if available, for straight quotes in verbatim environments
\IfFileExists{upquote.sty}{\usepackage{upquote}}{}
\IfFileExists{microtype.sty}{% use microtype if available
  \usepackage[]{microtype}
  \UseMicrotypeSet[protrusion]{basicmath} % disable protrusion for tt fonts
}{}
\usepackage{setspace}
\makeatletter
\@ifundefined{KOMAClassName}{% if non-KOMA class
  \IfFileExists{parskip.sty}{%
    \usepackage{parskip}
  }{% else
    \setlength{\parindent}{0pt}
    \setlength{\parskip}{6pt plus 2pt minus 1pt}}
}{% if KOMA class
  \KOMAoptions{parskip=half}}
\makeatother
\usepackage{longtable,booktabs,array}
\usepackage{caption}
\captionsetup[table]{skip=6pt}
\newcounter{none} % for unnumbered tables
\usepackage{calc} % for calculating minipage widths
% Correct order of tables after \paragraph or \subparagraph
\usepackage{etoolbox}
\makeatletter
\patchcmd\longtable{\par}{\if@noskipsec\mbox{}\fi\par}{}{}
\makeatother
% Allow footnotes in longtable head/foot
\IfFileExists{footnotehyper.sty}{\usepackage{footnotehyper}}{\usepackage{footnote}}
\makesavenoteenv{longtable}
\usepackage{graphicx}
\makeatletter
\newsavebox\pandoc@box
\newcommand*\pandocbounded[1]{% scales image to fit in text height/width
  \sbox\pandoc@box{#1}%
  \Gscale@div\@tempa{\textheight}{\dimexpr\ht\pandoc@box+\dp\pandoc@box\relax}%
  \Gscale@div\@tempb{\linewidth}{\wd\pandoc@box}%
  \ifdim\@tempb\p@<\@tempa\p@\let\@tempa\@tempb\fi% select the smaller of both
  \ifdim\@tempa\p@<\p@\scalebox{\@tempa}{\usebox\pandoc@box}%
  \else\usebox{\pandoc@box}%
  \fi%
}
% Set default figure placement to htbp
\def\fps@figure{htbp}
\makeatother
% definitions for citeproc citations
\NewDocumentCommand\citeproctext{}{}
\NewDocumentCommand\citeproc{mm}{%
  \begingroup\def\citeproctext{#2}\cite{#1}\endgroup}
\makeatletter
 % allow citations to break across lines
 \let\@cite@ofmt\@firstofone
 % avoid brackets around text for \cite:
 \def\@biblabel#1{}
 \def\@cite#1#2{{#1\if@tempswa , #2\fi}}
\makeatother
\newlength{\cslhangindent}
\setlength{\cslhangindent}{1.5em}
\newlength{\csllabelwidth}
\setlength{\csllabelwidth}{3em}
\newenvironment{CSLReferences}[2] % #1 hanging-indent, #2 entry-spacing
 {\begin{list}{}{%
  \setlength{\itemindent}{0pt}
  \setlength{\leftmargin}{0pt}
  \setlength{\parsep}{0pt}
  % turn on hanging indent if param 1 is 1
  \ifodd #1
   \setlength{\leftmargin}{\cslhangindent}
   \setlength{\itemindent}{-1\cslhangindent}
  \fi
  % set entry spacing
  \setlength{\itemsep}{#2\baselineskip}}}
 {\end{list}}
\usepackage{calc}
\newcommand{\CSLBlock}[1]{\hfill\break\parbox[t]{\linewidth}{\strut\ignorespaces#1\strut}}
\newcommand{\CSLLeftMargin}[1]{\parbox[t]{\csllabelwidth}{\strut#1\strut}}
\newcommand{\CSLRightInline}[1]{\parbox[t]{\linewidth - \csllabelwidth}{\strut#1\strut}}
\newcommand{\CSLIndent}[1]{\hspace{\cslhangindent}#1}
\setlength{\emergencystretch}{3em} % prevent overfull lines
\providecommand{\tightlist}{%
  \setlength{\itemsep}{0pt}\setlength{\parskip}{0pt}}
% Paragraph scaffolding for manuscript revision.
%
% Two environments, presented in this order per paragraph:
%
%   rgt   -- author's rewrite (initially a placeholder)
%            small, italic, dark blue
%   orig  -- original Claude-authored draft, and until the rewrite
%            replaces it, the body text of the manuscript
%            normal size, gray
%
% The `bullets' environment (a boiled-down topic sentence plus bullet
% points, one per paragraph) is no longer used inside manuscripts:
% those blocks now live in a companion bullets.Rmd supplement in the
% same report directory, which typesets them as plain \textbf{} +
% itemize and needs no part of this file. The definition is retained
% only so that a stray block left in a document still compiles.
%
% This file is identical across every manuscript that includes it.
% It previously existed in two variants that swapped the sizes of
% rgt and orig, so the same markup rendered differently depending on
% which repository the manuscript lived in. The variant retained here
% keeps orig at full size because orig, not rgt, currently carries
% the prose in every manuscript in the pool.

\usepackage{xcolor}

\newcounter{pgraph}

\newenvironment{bullets}{%
  \stepcounter{pgraph}%
  \begin{quote}\small\sffamily\color[gray]{0.30}%
  {\normalfont\scriptsize\color[gray]{0.58}[\thepgraph]\enspace}}{%
  \end{quote}}

\newenvironment{rgt}{%
  \begin{quote}\small\itshape\color[RGB]{20,60,160}}{%
  \end{quote}}

\newenvironment{orig}{%
  \begin{quote}\normalsize\color[gray]{0.55}}{%
  \end{quote}}

% backward compat alias
\newenvironment{claudecode}{\begin{orig}}{\end{orig}}
%% tools/stamp.tex
%% Provenance footer for PDFs rendered from this project.
%% Vendored by zzcollab; this file is not project-specific. The
%% footer prints the source document and the time of rendering.
%% The git version is defined here too, and so is available to a
%% project that wants it on the page, but it is left off the
%% default footer: it already names the staged copy in share/ and
%% appears in its manifest row. All values are supplied at render
%% time by tools/stamp-render.R, which writes a small generated
%% file that redefines the macros below. The \providecommand
%% defaults keep a bare LaTeX compile working when the document is
%% built without the wrapper.
\usepackage{fancyhdr}
\usepackage{xcolor}
\providecommand{\stampsource}{\detokenize{(source not recorded)}}
\providecommand{\stamptime}{(time not recorded)}
\providecommand{\stampversion}{\detokenize{(version not recorded)}}
\setlength{\footskip}{40pt}
\newcommand{\stampfootcontent}{%
  \footnotesize\color{gray}%
  \parbox[b]{\textwidth}{%
    \stamptime\hfill\thepage\\[2pt]%
    \ttfamily\raggedright\stampsource}}
\pagestyle{fancy}
\fancyhf{}
\fancyfoot[C]{\stampfootcontent}
\renewcommand{\headrulewidth}{0pt}
\renewcommand{\footrulewidth}{0.4pt}
%% The title page uses the `plain` page style; redefine it so the
%% stamp appears there too.
\fancypagestyle{plain}{%
  \fancyhf{}%
  \fancyfoot[C]{\stampfootcontent}%
  \renewcommand{\headrulewidth}{0pt}%
  \renewcommand{\footrulewidth}{0.4pt}}
\renewcommand{\stampsource}{\textasciitilde{}/\allowbreak{}prj/\allowbreak{}res/\allowbreak{}36-pmsimstats-ng/\allowbreak{}pmsimstats-ng/\allowbreak{}analysis/\allowbreak{}report/\allowbreak{}06-component-decomposition/\allowbreak{}report.Rmd}
\renewcommand{\stamptime}{2026-08-15 16:51 PDT}
\renewcommand{\stampversion}{\detokenize{bc0b671-wip}}
\usepackage{amsmath}
\usepackage{amssymb}
\usepackage{booktabs}
\usepackage{longtable}
\usepackage{float}
\usepackage[mathlines]{lineno}
\linenumbers
\usepackage{bookmark}
\IfFileExists{xurl.sty}{\usepackage{xurl}}{} % add URL line breaks if available
\urlstyle{same}
% fallback for those not using the hyperref driver hyperxmp:
\makeatletter
\@ifundefined{xmpquote}{\newcommand{\xmpquote}[1]{#1}}{}
\makeatother
\hypersetup{
  pdftitle={Three-component decomposition of treatment response in aggregated N-of-1 trials: pharmacological{,} expectation-driven{,} and natural-history components},
  pdfauthor={pmsimstats team},
  hidelinks,
  pdfcreator={LaTeX via pandoc}}

\title{Three-component decomposition of treatment response in aggregated N-of-1 trials: pharmacological, expectation-driven, and natural-history components}
\author{pmsimstats team}
\date{August 15, 2026}

\begin{document}
\maketitle

{
\setcounter{tocdepth}{2}
\tableofcontents
}
\setstretch{1.1}
\section*{Abstract}\label{abstract}
\addcontentsline{toc}{section}{Abstract}

\textbf{Background.} Treatment response in aggregated N-of-1 clinical
trials reflects three causally distinct mechanisms: pharmacological
action of the active compound (BR), the patient's belief that the
treatment will work (PB), and natural-history drift in the
underlying condition (TV). Standard analysis models combine these into
a single `treatment effect' plus residual error. Whether this
single-indicator combination biases inference is not unconditional: it depends on the
estimand and on whether candidate biomarkers correlate with the
non-pharmacological components, a dependence this paper makes
precise.

\textbf{Methods.} We formalize a three-component decomposition of the
within-participant outcome trajectory, each component a modified
Gompertz function with participant-specific random effects, and
characterize the trial-design contrasts (open-label, blinded
discontinuation, crossover) that supply identifiability. We derive
the omitted-variable-bias identity for the one-component estimator,
which expresses both the one-component treatment effect and the one-component
biomarker slope as sums of component-specific terms, and we
evaluate the analysis strategies in a pre-registered Monte Carlo
study (Study A) calibrated to the prazosin-PTSD reference
parameters, comparing a one-component design-contrast analysis with
a phase-augmented analysis at 1,000 replicates per alternative
cell (5,000 per null cell) across
\(m_{PB} \in \{0,1,3,6,10\}\), \(m_{TV} \in \{-1,0,1,2\}\), and
\(N \in \{35,70,200,300\}\).

\textbf{Results.} The identity shows that the one-component
biomarker-treatment slope is displaced from the pharmacological
slope by a term proportional to the biomarker's covariance with PB
and TV, not by the magnitude of those components. Study A confirms
this: with a biomarker coupled to BR only, the one-component
estimate of the biomarker-treatment interaction is essentially
unbiased across the whole grid (mean absolute bias \(0.016\) at
\(N=35\) and \(0.006\) at \(N=300\), within one Monte Carlo standard
error (MCSE) of zero), with near-nominal coverage and Type I error. The
phase-augmented analysis provides no bias reduction and is markedly
less powerful (mean power \(0.35\) versus \(0.59\) at \(N=35\); \(0.62\)
versus \(0.90\) at \(N=70\)), so it is dominated under this
uncontaminated biomarker. A contaminated-biomarker pilot then
exhibits the failure mode directly: coupling the biomarker to PB
drives the one-component bias up with the coupling strength,
reaching \(+0.51\) at a biomarker-PB correlation of \(0.45\) (\(N=300\),
some 30 MCSEs from zero), and the operative
quantity is the biomarker's correlation with the individual-level
PB \emph{variance}, not the population-mean placebo response, which
does not enter the slope; the phase-augmented analysis does not
remove this bias either. The inferential value of decomposition is
therefore conditional, not universal.

\textbf{Conclusions.} The trial design must supply the drug-on/off and
belief contrasts; given those, component-level decomposition is
useful when the estimand is the attribution of response to
pharmacology versus belief versus natural history, or when a
candidate biomarker may correlate with PB or TV. For a
biomarker-treatment interaction with a biomarker orthogonal to
PB and TV, a single-indicator analysis that exploits the design
contrast is sufficient and a parametric component model is
counterproductive. A pragmatic specification for the cases that
warrant it is
\texttt{Sx\ \textasciitilde{}\ bm\ +\ t\ +\ Dbc\ +\ phase\ +\ Dbc:phase\ +\ bm:Dbc\ +\ bm:Dbc:phase}
with random intercept and AR(1) residual correlation. A
contaminated-biomarker study demonstrates the failure mode (the
one-component slope is biased in proportion to the biomarker's correlation
with PB), and Study B shows that the bias is removable but only
under a design that decouples drug exposure from belief: on a
balanced-placebo design the belief-covariate decomposition recovers
the unbiased pharmacological slope (bias within one MCSE
of zero across the feasible contamination range at N=300) while the
one-component estimator does not, whereas on the standard hybrid
design no decomposition recovers it. Recovery is therefore
conditional on the trial design, not on the analysis alone.

\section{Introduction}\label{introduction}

\subsection{A claim about how N-of-1 trials should be analyzed}\label{a-claim-about-how-n-of-1-trials-should-be-analyzed}

\begin{rgt}
Lorem ipsum dolor sit amet, consectetur adipiscing elit, sed do eiusmod tempor incididunt ut labore et dolore magna aliqua.

\end{rgt}

\begin{orig}
This paper advances an opinionated methodological claim:
treatment response in aggregated N-of-1 clinical trials is the
sum of three causally distinct components and should routinely
be analyzed as such, rather than estimated as a single one-component
quantity. The three components are pharmacological action (\(BR\),
the biological response), placebo-belief response (\(PB\), the
patient's belief that they are receiving treatment), and
time-variant natural history (\(TV\), the trajectory that would
have occurred without intervention). Each is causally distinct,
each is independently of clinical and regulatory interest, and
each is identifiable from a properly designed N-of-1 trial. We
shall argue that the standard practice of fitting a single
mixed-effects model with one treatment indicator, prevalent
across the published N-of-1 literature, conflates these three
components, biases the interpretation of treatment effect, and
renders predictive-biomarker validation unanswerable
as a scientific question.

\end{orig}

\begin{rgt}
Lorem ipsum dolor sit amet, consectetur adipiscing elit, sed do eiusmod tempor incididunt ut labore et dolore magna aliqua.

\end{rgt}

\begin{orig}
We make no claim of novelty for the decomposition itself. The
structural decomposition of a longitudinal treatment response
into separable disease-progression, placebo, and drug-action
components is a well-established tradition in pharmacometrics\textsuperscript{\citeproc{ref-holford2015}{1}--\citeproc{ref-chen2021}{3}}, and the separation of
drug from placebo response in longitudinal trajectories has been
pursued both through growth-mixture models\textsuperscript{\citeproc{ref-muthen2009}{4}} and
through experimental decomposition of the placebo response
itself\textsuperscript{\citeproc{ref-kaptchuk2008components}{5}}. Each ingredient we use
(mixed-effects models with trial-phase indicators, placebo-arm
contrasts, regression-to-mean adjustment, growth-curve
parameterization) has likewise appeared across decades of
clinical-trial methodology. Our contribution is therefore one of
consolidation and extension rather than invention: we make the
three-component structure explicit for the aggregated N-of-1
setting, add a separately modeled placebo-belief channel
identified by design contrasts rather than by functional form
alone, and argue that the decomposition should become routine
wherever a longitudinal treatment response is analyzed. We
document why the framing has not been routine in N-of-1 practice
(parsimony, identifiability concerns, sample-size constraints,
methodological inertia) and propose an analysis that is
tractable at trial-relevant sample sizes through a
phase-augmented linear-mixed-effects approximation to the full
decomposition.

\end{orig}

\subsection{Why this matters: three failure modes of the one-component analysis}\label{why-this-matters-three-failure-modes-of-the-one-component-analysis}

\begin{rgt}
Lorem ipsum dolor sit amet, consectetur adipiscing elit, sed do eiusmod tempor incididunt ut labore et dolore magna aliqua.

\end{rgt}

\begin{orig}
We begin with the failure modes of the standard approach.
A one-component analysis (`the treatment effect is what we
observe minus baseline') has three identifiable failure modes
that arise when the population displays meaningful heterogeneity
in any of the unmodeled components. None of these is
hypothetical: each has a documented presence in the
chronic-condition trial literature. Two points of precision,
established formally in section 6.1 and borne out by the
simulation of section 7.3, should be carried through what
follows. The first two failure modes attach to the
baseline-anchored one-component effect and are removed once the design
supplies a drug-on/off contrast, which the one-component analysis
of later sections in fact uses. The third, biomarker-validation
failure, is different in kind: it requires not merely that \(PB\) or
\(TV\) be large but that the candidate biomarker correlate with
them, and where that coupling is absent the one-component analysis is not
biased at all.

\end{orig}

\begin{rgt}
Lorem ipsum dolor sit amet, consectetur adipiscing elit, sed do eiusmod tempor incididunt ut labore et dolore magna aliqua.

\end{rgt}

\begin{orig}
The first is the placebo-attribution failure. When the
population contains both high and low placebo responders, the
one-component treatment effect overestimates the pharmacological
component. The bias scales with the magnitude and variance of
\(PB\) in the study population, and it is invariant under
sample-size increases: a larger one-component trial converges
on the same biased estimate. The most consequential example in
modern clinical pharmacology is antidepressant efficacy.\textsuperscript{\citeproc{ref-kirsch2008}{6}} analyzed FDA-submitted data on four widely-
prescribed antidepressants and showed that, except at the most
severe baseline-severity stratum, mean drug-placebo differences
fell below clinical-significance thresholds, with the apparent
`efficacy' largely explained by placebo response and natural
fluctuation;\textsuperscript{\citeproc{ref-locher2015}{7}} reproduced this pattern in late-life
depression and explicitly invoked regression to the mean.\textsuperscript{\citeproc{ref-sugarman2014}{8}} generalized the finding to anxiety. Patient
populations differ systematically in placebo responsiveness,
and the difference has measurable biological correlates: COMT
val158met genotype predicts placebo response in irritable bowel
syndrome\textsuperscript{\citeproc{ref-hall2012comt}{9},\citeproc{ref-wang2022comt}{10}}, indicating that
`placebo' is not a homogeneous
nuisance but a structured biological phenomenon that interacts
with the patient's genome and treatment history. A trial whose
population happens to be placebo-responsive will report a
larger `drug effect' than one whose population is not, even
when the underlying pharmacology is identical.

\end{orig}

\begin{rgt}
Lorem ipsum dolor sit amet, consectetur adipiscing elit, sed do eiusmod tempor incididunt ut labore et dolore magna aliqua.

\end{rgt}

\begin{orig}
The second failure mode is regression-to-the-mean. When
patients enrol at moments of clinical engagement that coincide
with peak symptom severity, natural-history regression toward
typical severity contributes several points of apparent
improvement in the early weeks of any trial, without any
pharmacological contribution at all. The systematic reviews of\textsuperscript{\citeproc{ref-hrobjartsson2001}{11}},\textsuperscript{\citeproc{ref-hrobjartsson2004}{12}}, and\textsuperscript{\citeproc{ref-hrobjartsson2010}{13}}, the
last covering 202 trials across 60 clinical conditions,
quantified this directly: in trials with both placebo and
no-treatment arms, much of what is conventionally attributed to
the placebo response was indistinguishable from regression to
the mean and natural-history drift. The one-component analysis absorbs
this regression into the treatment-effect estimate and the
analyst has no formal mechanism to recover the true
pharmacological effect.

\end{orig}

\begin{rgt}
Lorem ipsum dolor sit amet, consectetur adipiscing elit, sed do eiusmod tempor incididunt ut labore et dolore magna aliqua.

\end{rgt}

\begin{orig}
The third failure mode is the biomarker-validation failure.
The precision-medicine question that motivates many recent
N-of-1 trials is whether a baseline biomarker predicts
treatment response, formalized in the PATH statement\textsuperscript{\citeproc{ref-kent2020path}{14},\citeproc{ref-kent2020pathee}{15}}. A biomarker that correlates
with placebo responsiveness or with a favorable natural-history
trajectory will appear in a one-component analysis as a
predictor of the one-component `treatment effect', indistinguishable
from a true pharmacological predictor. The biomarker-stratified
prescribing decision that follows from such an analysis is
wrong, and the trial's regulatory utility is degraded. The PATH
framework explicitly defers the mechanism of treatment-effect
heterogeneity to domain knowledge\textsuperscript{\citeproc{ref-kent2018bmj}{16},\citeproc{ref-rekkas2020scoping}{17}}; the BR-PB-TV decomposition supplies the
mechanistic substrate that PATH leaves blank. We argue that,
whenever the candidate biomarker may itself correlate with \(PB\) or
\(TV\), predictive-biomarker validation cannot be conducted
rigorously in a one-component analysis framework regardless of how large
the trial is; the bias is fixed by the biomarker's coupling to the
non-pharmacological components, not by sample size. Where the
biomarker can be assumed orthogonal to \(PB\) and \(TV\), this failure
mode does not arise, and the one-component analysis recovers the
pharmacological slope, as the simulation of section 7.3 confirms.
We stress at the outset the precise form this failure mode takes,
because the paper's own simulation history clarifies it. Study A
(section 7.3) varies the \emph{magnitude} of \(PB\) and \(TV\) with the
biomarker coupled to \(BR\) alone, which the section 6.1 identity
predicts to be innocuous, and it finds no bias. The decisive
experiment is the contaminated-biomarker pilot (section 7.4), which
couples the biomarker to \(PB\) and exhibits the failure directly: the
one-component bias grows with the biomarker-\(PB\) correlation. The operative
quantity is the biomarker's correlation with the individual-level
\(PB\) \emph{variance} (\(c_{bm,PB}\,\sigma_{PB}\)), not the
population-mean placebo response \(m_{PB}\), which does not enter the
slope; a trial population can have a large average placebo response
and yet pose no validation hazard, while one with zero average but a
biomarker-correlated spread of placebo responsiveness is at risk.
The failure-mode claims of this section are therefore consequences of
the identity, demonstrated in simulation (section 7.4), and the bias is
recoverable: Study B (section 7.5) shows that a belief-covariate
decomposition recovers the unbiased pharmacological slope under
contamination, but only on a design that decouples drug exposure from
belief, not on the standard hybrid design.

\end{orig}

\subsection{The current state of the N-of-1 literature}\label{the-current-state-of-the-n-of-1-literature}

\begin{rgt}
Lorem ipsum dolor sit amet, consectetur adipiscing elit, sed do eiusmod tempor incididunt ut labore et dolore magna aliqua.

\end{rgt}

\begin{orig}
The N-of-1 trial in its modern randomized form originates with\textsuperscript{\citeproc{ref-guyatt1986}{18}}, who defined the within-patient multiple-crossover
RCT, demonstrated it on a single airflow-limitation case, and
established a clinical service to deliver it at scale. The design
inherits its identifying logic from two older traditions: the
crossover trial of classical biostatistics, in which
within-subject period contrasts estimate treatment effects under
explicit carryover assumptions, and the single-case experimental
designs of behavioral science, in which phase reversals
(withdrawal and reinstatement of treatment) identify effects
within one participant. The blinded-discontinuation contrast
central to this paper is a direct descendant of both. The
methodological step from the individual-patient trial to
population-level inference was taken by\textsuperscript{\citeproc{ref-zucker1997}{19}} and\textsuperscript{\citeproc{ref-zucker2010}{20}}, who showed that a series of single-patient trials
can be combined through hierarchical meta-analysis to estimate
both population treatment effects and individual responses, the
statistical foundation on which the aggregated design rests. Over
the last fifteen years the design has accordingly been reframed
from an individual-patient decision tool into an aggregated,
meta-analytic framework supporting population-level inference:
the programmatic statement of\textsuperscript{\citeproc{ref-lillie2011}{21}}, the methodological
reviews of\textsuperscript{\citeproc{ref-duan2013}{22}}, the historical-practice audit of\textsuperscript{\citeproc{ref-gabler2011}{23}}, and the more recent syntheses of\textsuperscript{\citeproc{ref-schork2017nutrition}{24}},\textsuperscript{\citeproc{ref-schorkRandomizedClinicalTrials2018}{25}},\textsuperscript{\citeproc{ref-schork2023precision}{26}}, and\textsuperscript{\citeproc{ref-porcino2022}{27}} trace that trajectory. Reporting standards are
codified in the CONSORT
extension\textsuperscript{\citeproc{ref-vohraCONSORTExtensionReporting2015}{28}} (the headline statement) and\textsuperscript{\citeproc{ref-shamseer2015cent}{29}} (the explanation-and-elaboration companion);
the reporting-quality audits of\textsuperscript{\citeproc{ref-li2016cent}{30}} and its successors
document slow but progressive improvement in analytic
transparency. Three observations about this literature motivate
the present paper.

\end{orig}

\begin{rgt}
Lorem ipsum dolor sit amet, consectetur adipiscing elit, sed do eiusmod tempor incididunt ut labore et dolore magna aliqua.

\end{rgt}

\begin{orig}
The first observation is that the dominant analytic models are
simple. The reporting-quality audit of\textsuperscript{\citeproc{ref-li2016cent}{30}} documents that
the published N-of-1 trials between 1985 and 2013 relied
predominantly on paired t-tests, simple linear mixed-effects
models with a binary treatment indicator, or hierarchical
Bayesian regressions. On our own reading of that corpus and the
methodological literature that followed it, no published N-of-1
trial has employed a structural decomposition of treatment
response into pharmacological, expectancy, and natural-history
components. The most recent methodological
advances retain this character:\textsuperscript{\citeproc{ref-liao2023}{31}} develops Bayesian
distributed-lag models with autocorrelated errors as a current
carryover-aware analysis, but the model still treats the
response as a single quantity to be estimated.\textsuperscript{\citeproc{ref-blackstonComparisonAggregatedNof12019}{32}}
and\textsuperscript{\citeproc{ref-carrozzo2025}{33}} compare aggregated N-of-1 trials with
parallel and crossover RCT alternatives by simulation and
document carryover-induced Type I error inflation as a key
failure mode, but again at the level of the one-component treatment
effect. None of these papers decompose the response.

\end{orig}

\begin{rgt}
Lorem ipsum dolor sit amet, consectetur adipiscing elit, sed do eiusmod tempor incididunt ut labore et dolore magna aliqua.

\end{rgt}

\begin{orig}
The second observation is that the placebo-component separation
problem is widely known but largely unaddressed in N-of-1
specifically. The blinded discontinuation phase used in the
Hendrickson hybrid design\textsuperscript{\citeproc{ref-hendrickson2020}{34}} is one of the few
protocol-level mechanisms in routine N-of-1 use that supplies
identifiability of \(PB\) versus \(BR\). Other relevant approaches,
including placebo run-in designs, open-label-placebo trials, and
mediator-analysis applications, have been developed in parallel
but rarely imported into N-of-1 methodology.

\end{orig}

\begin{rgt}
Lorem ipsum dolor sit amet, consectetur adipiscing elit, sed do eiusmod tempor incididunt ut labore et dolore magna aliqua.

\end{rgt}

\begin{orig}
The third observation is that the closest precedent in the
literature is Kaptchuk's three-component decomposition of the
placebo response itself. In a randomized trial of
acupuncture-style placebo for irritable bowel syndrome,\textsuperscript{\citeproc{ref-kaptchuk2008components}{5}} decomposed the placebo response into
observation, ritual, and patient-practitioner relationship
components, demonstrating empirically that the placebo response
is itself a sum of identifiable parts rather than a single
`expectancy' lump. The follow-up open-label placebo trials of\textsuperscript{\citeproc{ref-kaptchuk2010olp}{35}},\textsuperscript{\citeproc{ref-lembo2021olp}{36}}, and\textsuperscript{\citeproc{ref-nurko2022}{37}} in adolescent
functional pain establish that even disclosed-placebo
intervention can produce clinically meaningful PB-only response.
The framework we propose is structurally analogous to
Kaptchuk's decomposition but applied at the level of the full
treatment response (drug, placebo, natural history) rather than
within the placebo response alone. To the best of our knowledge,
no published paper proposes the corresponding decomposition for
the aggregated N-of-1 setting, even though the methodological
ingredients, including mixed-effects modeling, blinded-phase
contrasts, and trajectory parameterization, are all in routine
use.

\end{orig}

\begin{rgt}
Lorem ipsum dolor sit amet, consectetur adipiscing elit, sed do eiusmod tempor incididunt ut labore et dolore magna aliqua.

\end{rgt}

\begin{orig}
The case we develop in the rest of this paper is therefore that
the gap is real, that the methodological tools to fill it are
mature, and that the failure modes of the one-component analysis are
large enough in chronic-condition pharmacology to make the
decomposition a methodological priority rather than an academic
refinement.

\end{orig}

\subsection{The neurobiological and epidemiological reality of the components}\label{the-neurobiological-and-epidemiological-reality-of-the-components}

\begin{rgt}
Lorem ipsum dolor sit amet, consectetur adipiscing elit, sed do eiusmod tempor incididunt ut labore et dolore magna aliqua.

\end{rgt}

\begin{orig}
The decomposition is not a statistical convenience. Each
component reflects a real causal mechanism documented in
independent literatures, and we shall briefly recall each in
turn before turning to the formal model.

\end{orig}

\begin{rgt}
Lorem ipsum dolor sit amet, consectetur adipiscing elit, sed do eiusmod tempor incididunt ut labore et dolore magna aliqua.

\end{rgt}

\begin{orig}
The biological response component represents the pharmacological
action of the active compound: the chemical binding, downstream
signaling, and physiological response that a regulatory
authority cares about when it grants approval. The
placebo-belief response reflects neurobiologically real
mechanisms documented across pain, depression, Parkinson's, and
PTSD research\textsuperscript{\citeproc{ref-benedetti2014}{38}--\citeproc{ref-carlino2016benedetti}{41}}, mediated by belief-driven release of
endogenous opioids and dopamine, anticipatory regulation of the
hypothalamic-pituitary-adrenal axis, and top-down attentional
shifts in symptom perception. Direct neuroimaging evidence is
instructive here:\textsuperscript{\citeproc{ref-wager2013nps}{42}} developed an fMRI-based
`neurologic pain signature' (NPS) that tracks nociceptive input
and is reduced by the opioid agonist remifentanil, and the
individual-participant meta-analysis of\textsuperscript{\citeproc{ref-zunhammer2018nps}{43}} found
that placebo treatments produce moderate reductions in reported pain
but only negligible reductions in the NPS, providing a neural
dissociation between the pharmacological and expectancy
contributions to a clinical pain response. The placebo response
is heterogeneous across patients with biological correlates\textsuperscript{\citeproc{ref-hall2012comt}{9},\citeproc{ref-wang2022comt}{10}} and is
modulated by trial phase: when belief varies, as it does
between open-label and blinded-discontinuation conditions, the
placebo response varies, and the modulation is exactly the
design contrast that makes the decomposition identifiable.

\end{orig}

\begin{rgt}
Lorem ipsum dolor sit amet, consectetur adipiscing elit, sed do eiusmod tempor incididunt ut labore et dolore magna aliqua.

\end{rgt}

\begin{orig}
The time-variant natural-history component reflects whatever
trajectory the participant would have followed without treatment.
For acute conditions this trends positive (recovery on intrinsic
timescale); for chronic stable conditions it is near zero on
the population mean but with substantial participant-level
variance; for chronic progressive conditions it trends negative;
and for selection-driven enrollment processes it is positive in
the early weeks of the trial regardless of underlying disease
class, because of regression to the mean\textsuperscript{\citeproc{ref-hrobjartsson2001}{11},\citeproc{ref-hrobjartsson2010}{13}}. The PATH statement\textsuperscript{\citeproc{ref-kent2020path}{14},\citeproc{ref-kent2020pathee}{15}} highlights that this
heterogeneity is itself the central object of precision-medicine
inference: HTE must be modeled,
not absorbed into noise, if the precision-medicine question is
to be answered.

\end{orig}

\subsection{What this paper contributes}\label{what-this-paper-contributes}

\begin{rgt}
Lorem ipsum dolor sit amet, consectetur adipiscing elit, sed do eiusmod tempor incididunt ut labore et dolore magna aliqua.

\end{rgt}

\begin{orig}
We make four contributions, each addressed in a section below.
We begin in section 2 with the formal model and notation.
Section 3 develops each component in turn, with biological
motivation, parametric form, and identifiability conditions.
Section 4 maps the design-contrast structure of the standard
hybrid open-label-plus-blinded-discontinuation design onto the
identifiability requirements of the decomposition. Section 5
develops the analysis-side methodology, including the
linear-mixed-effects approximations and the phase-by-treatment
indicator extension that retains BR-PB sensitivity at low
parametric cost. Section 6 presents the worked examples, and
section 7 presents the simulation study. Section 8 discusses
applications to predictive-biomarker validation. The companion
pedagogical document at
\texttt{docs/24-component-decomposition-pedagogy.md} provides an
extended worked-example treatment for readers approaching the
material for the first time.

\end{orig}

\section{Notation and the formal model}\label{notation-and-the-formal-model}

\begin{rgt}
Lorem ipsum dolor sit amet, consectetur adipiscing elit, sed do eiusmod tempor incididunt ut labore et dolore magna aliqua.

\end{rgt}

\begin{orig}
In this section we shall fix notation and state the formal
model. The three components are defined precisely enough to make
the identifiability and analysis-side arguments in later
sections rigorous, but we have endeavored to keep the
presentation compact. A fuller treatment, with worked numerical
examples at each step, is available in the companion pedagogical
document at \texttt{docs/24-component-decomposition-pedagogy.md}.

\end{orig}

\begin{rgt}
Lorem ipsum dolor sit amet, consectetur adipiscing elit, sed do eiusmod tempor incididunt ut labore et dolore magna aliqua.

\end{rgt}

\begin{orig}
For participant \(i\) at trial timepoint \(t\):

\[
Y_{it} \;=\; \mathrm{BL}_i \;-\; \bigl[\,BR_{it} + PB_{it} +
TV_{it}\,\bigr] \;+\; \varepsilon_{it},
\]

\end{orig}

\begin{rgt}
Lorem ipsum dolor sit amet, consectetur adipiscing elit, sed do eiusmod tempor incididunt ut labore et dolore magna aliqua.

\end{rgt}

\begin{orig}
Each component is a participant-specific function of its own
time-on-state variable. \(BR_{it}\) depends on time-on-drug, the
cumulative exposure to active medication since drug initiation.
\(PB_{it}\) depends on time-on-belief, modulated by a
phase-specific expectancy factor \(\eta(\text{phase})\) that
captures the patient's confidence that they are receiving active
treatment. \(TV_{it}\) depends on time-since-trial-entry,
calendar time elapsed regardless of treatment status.

\end{orig}

\begin{rgt}
Lorem ipsum dolor sit amet, consectetur adipiscing elit, sed do eiusmod tempor incididunt ut labore et dolore magna aliqua.

\end{rgt}

\begin{orig}
Each component follows a modified Gompertz curve:

\[
C(t) \;=\; m_C \,\exp\bigl(-d_C \exp(-r_C t)\bigr) \,-\,
m_C\,\exp(-d_C),
\]

\end{orig}

\begin{rgt}
Lorem ipsum dolor sit amet, consectetur adipiscing elit, sed do eiusmod tempor incididunt ut labore et dolore magna aliqua.

\end{rgt}

\begin{orig}
\(\varepsilon_{it}\) is a within-participant residual term with
AR(1) serial-correlation structure
\(\mathrm{Cov}(\varepsilon_{it}, \varepsilon_{i,t+h}) =
\sigma^2_\varepsilon \rho^{|h|}\). The three components are
permitted to be jointly correlated through within-time and
across-time cross-component covariances \(c_{cf1t}\) and
\(c_{cfct}\) respectively, so that participant-level heterogeneity
in any one component can co-vary with heterogeneity in the
others.

\end{orig}

The principal psychometric and pharmacological terms used below
are summarized in Table \ref{tab:terms}.

\begin{table}

\caption{\label{tab:terms}Principal psychometric and pharmacological terms used in the three-component decomposition.}
\centering
\begin{tabular}[t]{ll}
\toprule
Term & Meaning\\
\midrule
BR & Biological response: pharmacological reduction in symptom score attributable to active drug.\\
PB & Placebo-belief response: reduction attributable to the patient's belief that they are receiving treatment.\\
TV & Time-variant component: reduction (or increase) attributable to the natural-history trajectory of the underlying condition.\\
$\eta(\text{phase})$ & Expectancy multiplier for $PB$: 1 in open-label on-drug, $\approx 0.5$ in blinded discontinuation, between $0$ and $0.5$ in open-label placebo.\\
$\mathrm{BL}_i$ & Participant $i$'s baseline (pre-treatment) symptom score.\\
\addlinespace
$D_{it}$ & Continuous-decay drug indicator on the analysis side: 1 on-drug, $\exp(-\lambda t_{sd,it})$ off-drug.\\
$t_{sd,it}$ & Time since drug discontinuation (off-drug only).\\
$\lambda$ & Carryover decay rate, with half-life $t_{1/2} = \ln 2 / \lambda$.\\
\bottomrule
\end{tabular}
\end{table}

\section{The three components}\label{the-three-components}

\begin{rgt}
Lorem ipsum dolor sit amet, consectetur adipiscing elit, sed do eiusmod tempor incididunt ut labore et dolore magna aliqua.

\end{rgt}

\begin{orig}
In this section we shall present the three components in turn,
each with its biological motivation, the modified Gompertz
parameterization, and the design contrasts that identify it.
The treatment is organized so that the reader can see, by the
end of this section, that the formal model is a sum of three
causally interpretable trajectories rather than a mathematical
convenience.

\end{orig}

\subsection{Biological response (BR)}\label{biological-response-br}

\begin{rgt}
Lorem ipsum dolor sit amet, consectetur adipiscing elit, sed do eiusmod tempor incididunt ut labore et dolore magna aliqua.

\end{rgt}

\begin{orig}
\(BR\) is the physiological reduction in symptom score attributable
to the active medication's chemical action on the body. It is
independent of belief and natural history: a
participant who somehow received active drug without knowing it
would still develop a measurable \(BR\), and a participant who
took an inert sugar pill while believing it was active drug
would have \(BR = 0\) no matter how strong their expectation.
Methodologically, \(BR\) is the quantity that drug regulators and
prescribers want: approval decisions, dose selection, and
biomarker-guided patient stratification are all framed around
\(BR\).

\end{orig}

\begin{rgt}
Lorem ipsum dolor sit amet, consectetur adipiscing elit, sed do eiusmod tempor incididunt ut labore et dolore magna aliqua.

\end{rgt}

\begin{orig}
The qualitative shape of the \(BR\) trajectory follows from
pharmacokinetic and pharmacodynamic considerations. Drug
concentration in plasma reaches steady state within a small
multiple of the elimination half-life; brain concentration
follows plasma with a short delay; the downstream
pharmacodynamic response (reduction in symptom-relevant neural
activity, in the case of prazosin's effect on noradrenergic tone
during REM sleep) typically requires several days to weeks of
sustained exposure to integrate fully\textsuperscript{\citeproc{ref-raskind2013}{44}}.
The result is a saturating curve: slow rise in the first week,
steepest growth in the second to third week, plateau as receptor
adaptation completes, and a decay back toward zero on
discontinuation governed by the carryover half-life
\(t_{1/2} = \ln 2 / \lambda\) rather than an instantaneous drop.
The rate of this off-drug decay is the object of the carryover
machinery used throughout this project; only in the limit of a
short half-life relative to the measurement spacing does it
approximate the sharp step assumed in the worked example of
section 6.

\end{orig}

\begin{rgt}
Lorem ipsum dolor sit amet, consectetur adipiscing elit, sed do eiusmod tempor incididunt ut labore et dolore magna aliqua.

\end{rgt}

\begin{orig}
The modified Gompertz form
\(BR(t) = m_{BR}\exp(-d_{BR}\exp(-r_{BR}t)) -
m_{BR}\exp(-d_{BR})\) captures this profile compactly. The
asymptote \(m_{BR}(1 - e^{-d_{BR}})\) is the biological ceiling,
the largest reduction the drug can produce at infinite exposure;
the scale parameter \(m_{BR}\) equals that ceiling up to the
factor \((1 - e^{-d_{BR}})\), which for the calibrated
\(d_{BR} \approx 5\) exceeds \(0.99\). The onset
rate \(r_{BR}\) controls how rapidly the participant approaches the
ceiling: high \(r_{BR}\) means a fast onset (effect plateaus in
one to two weeks), low \(r_{BR}\) a slow climb (six to eight
weeks). The displacement \(d_{BR}\) shapes the climb: high
\(d_{BR}\) gives an early-onset curve that rises sharply at first,
low \(d_{BR}\) a more gradual climb. For PTSD prazosin, calibrated
values from the\textsuperscript{\citeproc{ref-hendrickson2020}{34}} reference dataset are
\(m_{BR} \approx 11\) (eleven points of nightmare-score reduction
at saturation), \(r_{BR} \approx 0.42\) per week (half-maximum at
roughly four to five weeks of sustained exposure), and
\(d_{BR} \approx 5\) (moderately steep early climb).

\end{orig}

\begin{rgt}
Lorem ipsum dolor sit amet, consectetur adipiscing elit, sed do eiusmod tempor incididunt ut labore et dolore magna aliqua.

\end{rgt}

\begin{orig}
Two alternative families that do not capture the asymmetric
saturation profile are worth noting for contrast. A linear
function \(BR = \beta t\) climbs without bound and is biologically
implausible past a few weeks. An exponential decay
\(BR = m_{BR}(1 - \exp(-r t))\) is acceptable for fast-onset drugs
but lacks the slow-start phase that characterizes most
chronic-condition responses, where receptor adaptation, feedback
regulation, and gene-expression changes produce a measurable
lag. The Gompertz interpolates between exponential decay (high
\(d\) limit) and linear-in-\(t\) growth (low \(d\) limit), and its
three-parameter form is flexible enough to fit a wide range of
pharmacodynamic profiles without overfitting. The Gompertz is a
deliberate choice over the two parametric forms most familiar
from pharmacokinetic-pharmacodynamic modeling. The sigmoid
\(E_{\max}\) (Hill) model\textsuperscript{\citeproc{ref-holford2011backstory}{45}} describes the
dependence of effect on concentration (a dose-response curve),
whereas \(BR(t)\) here is a time-course of accumulating effect at
fixed dosing; and the indirect-response (turnover) models that
do describe effect time-courses are specified through
differential equations on a hidden response variable rather than
in closed form. The modified Gompertz retains the slow-start,
saturating shape those models produce while remaining a
closed-form, three-parameter mean function that is tractable
inside a mixed-effects likelihood. Doc 25 of the
project documentation provides the full historical and
biological context for the family choice and quantifies the
sensitivity of pmsimstats inference to alternative families;
that work is reported separately as \texttt{07-gompertz-evaluation}.
The headline result of that evaluation is reassuring for the present
framework: in a 1\{,\}000-replicate-per-cell simulation that generated
data under the logistic, hyperbolic-tangent, and piecewise-linear
breakpoint families and analyzed them under the standard linear-mixed
model, the trajectory family was essentially irrelevant to the power of
the biomarker-treatment interaction test (the family-to-family spread
fell within two MCSEs of sampling noise) and to
Type I error (0.025 to 0.036 across null cells, with no family-specific
inflation). The conclusions drawn here are therefore not artifacts of
the Gompertz form, though the regimes the family cannot represent
(cyclical or biphasic trajectories) remain genuine exclusions.

\end{orig}

\begin{rgt}
Lorem ipsum dolor sit amet, consectetur adipiscing elit, sed do eiusmod tempor incididunt ut labore et dolore magna aliqua.

\end{rgt}

\begin{orig}
The \(BR\) component is identified by within-participant on-drug
versus off-drug contrasts, formalized in section 4 below. The
most direct contrast is the blinded-discontinuation phase: when
a participant is silently switched to placebo while still
believing they are on active drug, \(BR\) decays toward zero as the
drug clears while \(PB\) remains roughly intact, and the difference
between on-drug and off-drug response in that window is a
relatively clean estimate of \(BR\) that does not require the
analyst to model \(PB\) or \(TV\) explicitly. The decay is not
instantaneous: with the carryover half-life
\(t_{1/2} = \ln 2 / \lambda\) of the analysis-side decay term,
residual \(BR\) persists into the early off-drug observations, so
the contrast is clean only once the discontinuation window is
long relative to \(t_{1/2}\), and the drug state is best carried by
the continuous-decay indicator \(D_{it}\) of Table
\ref{tab:terms} rather than a hard on/off indicator.

\end{orig}

\subsection{Placebo-belief response (PB)}\label{placebo-belief-response-pb}

\begin{rgt}
Lorem ipsum dolor sit amet, consectetur adipiscing elit, sed do eiusmod tempor incididunt ut labore et dolore magna aliqua.

\end{rgt}

\begin{orig}
\(PB\) is the reduction in symptom score that arises from the
patient's belief that they are receiving active treatment,
independent of what they are actually receiving. That is, it is
the placebo response in its strict sense: the part of the
improvement that would persist if the active drug were silently
replaced with an inert pill, provided the patient continued to
believe they were on the active drug.

\end{orig}

\begin{rgt}
Lorem ipsum dolor sit amet, consectetur adipiscing elit, sed do eiusmod tempor incididunt ut labore et dolore magna aliqua.

\end{rgt}

\begin{orig}
It is worth being precise about what \(PB\) is and is not. \(PB\)
is not `fake' improvement; it is a phenomenon clinicians
knowingly exploit, a national survey of US internists and
rheumatologists finding roughly half report regular use of
placebo treatments\textsuperscript{\citeproc{ref-tilburt2008}{46}}. As the introduction sets out,
the placebo response is mediated by well-characterized
neurobiological pathways\textsuperscript{\citeproc{ref-benedetti2014}{38}--\citeproc{ref-carlino2016benedetti}{41},\citeproc{ref-finniss2010}{47}}:
belief-driven release of endogenous opioids and dopamine,
anticipatory hypothalamic-pituitary-adrenal regulation, and
top-down attentional shifts, dissociated from pharmacology by the\textsuperscript{\citeproc{ref-wager2013nps}{42}} neurologic pain signature\textsuperscript{\citeproc{ref-zunhammer2018nps}{43}} and
shown to have genetic correlates\textsuperscript{\citeproc{ref-hall2012comt}{9},\citeproc{ref-wang2022comt}{10}}.
What makes these mechanisms `placebo' rather than `drug' is that
they are triggered by belief rather than by pharmacology.

\end{orig}

\begin{rgt}
Lorem ipsum dolor sit amet, consectetur adipiscing elit, sed do eiusmod tempor incididunt ut labore et dolore magna aliqua.

\end{rgt}

\begin{orig}
For trial design, the relevant point is that \(PB\) is separately
identifiable from \(BR\) if and only if the trial includes a phase
in which the patient's belief about treatment differs from the
actual treatment they are receiving. The blinded-discontinuation
phase is the canonical such phase: the patient is told that, at
some unannounced point during the window, they may be silently
switched to placebo, and they cannot reliably tell whether they
are still receiving active drug. Belief in treatment
correspondingly drops to a partial level (multiplier roughly
\(0.5\)) because the patient's expectation is hedged.

\end{orig}

\begin{rgt}
Lorem ipsum dolor sit amet, consectetur adipiscing elit, sed do eiusmod tempor incididunt ut labore et dolore magna aliqua.

\end{rgt}

\begin{orig}
\(PB\) also follows a Gompertz curve, scaled by the expectancy
factor:

\[
PB(t \mid \text{phase}) \;=\; \eta(\text{phase}) \,\cdot\,
m_{PB}\,\exp\bigl(-d_{PB}\exp(-r_{PB} t)\bigr)
\,+\, (\text{offset to start at zero}).
\]

Typical expectancy values are \(\eta = 1.0\) in open-label
on-drug, \(\eta \approx 0.5\) in blinded discontinuation, and
\(\eta\) between \(0\) and \(0.5\) in open-label placebo (the residual
expectation a patient retains when told the previous active
phase was beneficial). The quantity \(\eta\) operationalizes
response expectancy in the sense of\textsuperscript{\citeproc{ref-kirsch1985}{48}}, the patient's
anticipation of benefit that the expectancy literature treats as
a determinant of experienced symptom change. The choice of
\(\eta = 0.5\) for the blinded phase is a modeling assumption
rather than a calibrated quantity: to our knowledge the residual expectancy a patient
carries when uncertain about allocation has not been measured
directly as a fraction of the open-label level, and we adopt
\(0.5\) as a deliberately central placeholder. The available
indirect evidence is nonetheless that the residual response is
substantial. Open-label placebo trials\textsuperscript{\citeproc{ref-kaptchuk2010olp}{35}--\citeproc{ref-nurko2022}{37}} demonstrate that a clinically
meaningful response persists even when patients are explicitly
told they are receiving an inert intervention, which bounds the
residual belief-driven response away from zero. Because the
value of \(\eta\) is not empirically pinned, sensitivity analyses
varying it across its plausible range are mandatory before any
conclusion is allowed to depend on this single number.

\end{orig}

\begin{rgt}
Lorem ipsum dolor sit amet, consectetur adipiscing elit, sed do eiusmod tempor incididunt ut labore et dolore magna aliqua.

\end{rgt}

\begin{orig}
The decomposition framework as proposed here has a structural
analogue in the placebo-research literature.\textsuperscript{\citeproc{ref-kaptchuk2008components}{5}}
decomposed the placebo response itself into observation, ritual,
and patient-practitioner relationship components in a randomized
trial of acupuncture-style placebo for irritable bowel syndrome,
demonstrating empirically that the placebo response is not a
single `expectancy' lump but a sum of identifiable parts. The
present BR-PB-TV framework is analogous in spirit but operates
at the level of the full treatment response (drug, placebo,
natural history) rather than within the placebo response alone.

\end{orig}

\begin{rgt}
Lorem ipsum dolor sit amet, consectetur adipiscing elit, sed do eiusmod tempor incididunt ut labore et dolore magna aliqua.

\end{rgt}

\begin{orig}
A subtle property of \(PB\) is that its variance also scales with
\(\eta\), not just its mean. When \(\eta\) is high, there is
substantial variability across patients in how strongly the
placebo response is expressed; high-expectancy patients show
large \(PB\), low-expectancy patients show essentially none.
When \(\eta\) is low (blinded), the placebo response is uniformly
attenuated and the between-participant variance shrinks
accordingly. The model therefore parameterizes
\(\mathrm{SD}(PB \mid \text{phase}) = \eta(\text{phase}) \cdot
\sigma_{PB}^{\text{baseline}}\). Ignoring this heteroscedastic
structure deflates standard errors in the open-label phase and
inflates them in the blinded phase, with downstream consequences
for confidence-interval coverage and for the calibration of any
biomarker test built on top of the analysis.

\end{orig}

\subsubsection{Additivity, subadditivity, and the BR-PB interaction}\label{additivity-subadditivity-and-the-br-pb-interaction}

\begin{rgt}
Lorem ipsum dolor sit amet, consectetur adipiscing elit, sed do eiusmod tempor incididunt ut labore et dolore magna aliqua.

\end{rgt}

\begin{orig}
In this section we shall broach a question that the formal
model leaves open: whether the three components contribute
additively to symptom reduction, or whether there is a
meaningful interaction between \(BR\) and \(PB\) that the additive
specification misses.

\end{orig}

\begin{rgt}
Lorem ipsum dolor sit amet, consectetur adipiscing elit, sed do eiusmod tempor incididunt ut labore et dolore magna aliqua.

\end{rgt}

\begin{orig}
The decomposition as written in section 2 treats the three
components as additive contributors to symptom reduction:

\[
Y_{it} \;=\; \mathrm{BL}_i \;-\; (BR_{it} + PB_{it} + TV_{it})
            \;+\; \varepsilon_{it}.
\]

The additivity is a modeling assumption, not a fact about
biology, and it is worth stating that the assumption may be
wrong in directions the framework can absorb and in directions
it cannot. Additivity is also scale-dependent in the sense
emphasized by\textsuperscript{\citeproc{ref-senn2004}{49}}: components that combine
non-additively on the raw symptom scale may combine additively
on a transformed (for example logarithmic) scale, so an apparent
\(BR \times PB\) interaction can be partly an artifact of the
chosen response metric rather than a biological fact. That
belief and pharmacology interact rather than act in isolation is
demonstrated directly by\textsuperscript{\citeproc{ref-bingel2011}{50}}, who showed with
neuroimaging that positive treatment expectation roughly doubled
the analgesic benefit of the opioid remifentanil while negative
expectation abolished it, so the two channels are not in general
separable by simple addition. More specifically, the published
literature on placebo-drug interactions
has accumulated evidence that the combined effect of receiving
an active drug and believing one is receiving an active drug is,
in some clinical contexts, less than the sum of the two
components measured separately\textsuperscript{\citeproc{ref-kaptchuk2008components}{5},\citeproc{ref-benedetti2014}{38}}. The phenomenon is typically termed
subadditivity: \(\partial Y / \partial(BR \cdot PB) > 0\) when
both components are present and their interaction term is
included in a saturated mean-structure model.

\end{orig}

\begin{rgt}
Lorem ipsum dolor sit amet, consectetur adipiscing elit, sed do eiusmod tempor incididunt ut labore et dolore magna aliqua.

\end{rgt}

\begin{orig}
Three mechanisms have been proposed in the placebo-research
literature to explain subadditivity. First, a shared neural
substrate: if the descending pain-modulation circuits that
mediate placebo analgesia overlap with the circuits modulated
by the active drug, simultaneous activation by both mechanisms
produces less marginal benefit than activation by either alone,
in a saturating-occupancy sense. Second, ceiling effects on the
symptom scale: a clinical scale that caps at zero (no symptoms)
cannot register additional reduction once a participant has
reached that cap, so any participant whose \(BR + PB\) already
saturates the scale will register no further reduction even if
\(BR\) or \(PB\) would, in isolation, predict more. Third,
belief-dynamics adaptation: the patient's expectation may adapt
downward when the active drug produces a strong response that
exceeds the expected effect, reducing the ongoing \(PB\)
contribution; equivalently, unexpectedly large \(BR\) may attenuate
\(PB\) via expectation recalibration.

\end{orig}

\begin{rgt}
Lorem ipsum dolor sit amet, consectetur adipiscing elit, sed do eiusmod tempor incididunt ut labore et dolore magna aliqua.

\end{rgt}

\begin{orig}
Whichever mechanism dominates in a given clinical context, the
subadditivity question is empirical and design-dependent.
Within the BR-PB-TV framework, subadditivity manifests as a
non-zero coefficient on a \(BR \times PB\) interaction term in
the conditional-mean specification. The additive form treats
this interaction as zero by construction; a generalized form

\[
Y_{it} \;=\; \mathrm{BL}_i \;-\; \bigl[BR_{it} + PB_{it}
            + TV_{it} + \gamma \cdot BR_{it} \cdot PB_{it}\bigr]
            \;+\; \varepsilon_{it},
\]

\end{orig}

\begin{rgt}
Lorem ipsum dolor sit amet, consectetur adipiscing elit, sed do eiusmod tempor incididunt ut labore et dolore magna aliqua.

\end{rgt}

\begin{orig}
For trial design, two consequences follow. First, a biomarker
that predicts \(BR\) in an additive-decomposition analysis may
overstate or understate predictive accuracy if the true
mechanism is subadditive: under subadditivity, the
biomarker-driven increase in \(BR\) is partly cancelled by the
\(\gamma \cdot BR \cdot PB\) interaction, and the marginal
biomarker-treatment association is smaller than the
biomarker-by-\(BR\) association. Second, trial designs that
attenuate \(PB\) in some phases (blinded discontinuation,
crossover) implicitly down-weight the subadditivity contribution
to the response in those phases, which means that the effective
biomarker-treatment association estimated from those phases is
closer to the biomarker-by-\(BR\) association than is the
open-label estimate. The phase-augmented analysis recommended
in section 5 below therefore has the side benefit of partly
absorbing subadditivity, even when \(\gamma\) is not estimated
explicitly.

\end{orig}

\begin{rgt}
Lorem ipsum dolor sit amet, consectetur adipiscing elit, sed do eiusmod tempor incididunt ut labore et dolore magna aliqua.

\end{rgt}

\begin{orig}
A simulation study quantifying the magnitude of bias in the
biomarker-treatment interaction estimate when the additive
analysis is fitted to data generated under non-zero \(\gamma\) is
forthcoming as part of the simulation program described in
section 7. The pre-registration is that bias scales linearly
with \(\gamma\) at small \(|\gamma|\), with the sign opposite to the
true biomarker-by-\(BR\) effect under subadditivity and aligned
with it under superadditivity.

\end{orig}

\begin{rgt}
Lorem ipsum dolor sit amet, consectetur adipiscing elit, sed do eiusmod tempor incididunt ut labore et dolore magna aliqua.

\end{rgt}

\begin{orig}
A subtle but inferentially important point merits attention
here: a non-zero \(\hat{\gamma}\) recovered from an
additive-extension analysis does not, on its own, establish
that the underlying mechanism is subadditive interaction in the
strict sense. The same marginal pattern can arise from latent
population structure that the analyst is not modeling. If the
trial population is a mixture of treatment-responders and
non-responders (in the latent-class sense developed at length
in \texttt{analysis/report/03-latent-class-mixture-application/}),
and class membership correlates with both \(BR\) strength and
\(PB\) strength (responder-class participants have both higher
\(BR\) and higher \(PB\) on average), then the marginal expected
value of \(BR \cdot PB\) over the unobserved class label has a
non-zero covariance term whose magnitude depends on the gating
function \(\pi(B)\). From the one-component analysis perspective this
looks exactly like a \(\gamma\)-weighted \(BR \times PB\) interaction
even though no such direct interaction was specified in the
class-conditional DGP. That is, the latent-class structure is
one mechanism that produces apparent subadditivity in a
one-component analysis, regardless of whether the underlying
biology has a true direct \(BR\)-\(PB\) interaction.

\end{orig}

\begin{rgt}
Lorem ipsum dolor sit amet, consectetur adipiscing elit, sed do eiusmod tempor incididunt ut labore et dolore magna aliqua.

\end{rgt}

\begin{orig}
The two mechanisms, namely true subadditive interaction and
latent-class structure with class-correlated components,
produce qualitatively similar marginal effects but are
inferentially distinct. Distinguishing them requires either a
class-aware analysis that explicitly models the latent
population structure (as the latent-class paper develops at
length), or a phase contrast that separates the mechanisms by
exploiting the differential expectancy modulation of \(PB\) that
the BR-PB-TV framework provides. Under the additive-extension
analysis fitted to mixture-DGP data with class-correlated
components, \(\hat{\gamma}\) is identified up to confounding with
the class structure, and the analyst should report
\(\hat{\gamma}\) alongside a class-aware sensitivity analysis
rather than as standalone evidence for mechanistic
subadditivity. The Study 1 simulation in the latent-class paper
includes a \texttt{pb\_class\_correlation} axis that varies
\(\mathrm{Cor}(BR, PB)\) across class labels; that axis is the
direct simulation-side test of the latent-class generation of
apparent subadditivity, and the cross-paper synthesis between
the two papers is the cleanest way to disentangle the two
mechanisms.

\end{orig}

\subsection{Time-variant natural-history component (TV)}\label{time-variant-natural-history-component-tv}

\begin{rgt}
Lorem ipsum dolor sit amet, consectetur adipiscing elit, sed do eiusmod tempor incididunt ut labore et dolore magna aliqua.

\end{rgt}

\begin{orig}
\(TV\) is the reduction in symptom score that would have happened
even if the patient had received no treatment at all. It is the
participant's natural-history trajectory, integrated over the
duration of the trial. For some participants \(TV\) is positive
(symptoms naturally improve over the trial because of life
events, therapy, time-in-condition, or regression to the mean);
for others \(TV\) is negative (symptoms naturally worsen because
of trauma anniversaries, accumulating stress, or progressive
underlying disease).

\end{orig}

\begin{rgt}
Lorem ipsum dolor sit amet, consectetur adipiscing elit, sed do eiusmod tempor incididunt ut labore et dolore magna aliqua.

\end{rgt}

\begin{orig}
Like \(PB\), \(TV\) is not noise. It is a structured
participant-specific signal that the model can in principle
estimate from data, given sufficient repeated measurements and
a trial design that includes a long enough off-drug baseline or
follow-up window. What distinguishes \(TV\) from \(PB\) is its
origin: \(TV\) is independent of treatment status and of belief
about treatment, whereas \(PB\) is modulated by both.

\end{orig}

\subsubsection{When does natural history actually exist?}\label{when-does-natural-history-actually-exist}

\begin{rgt}
Lorem ipsum dolor sit amet, consectetur adipiscing elit, sed do eiusmod tempor incididunt ut labore et dolore magna aliqua.

\end{rgt}

\begin{orig}
A common simplifying assumption in trial analysis is that
subjects will improve over time, with or without treatment.
This is true for some conditions and problematic for others.
The shape of \(TV\) depends on the underlying disease class, the
trial timescale, and the selection process by which participants
entered the trial. None of these can be assumed without thought,
and the empirical shape of \(TV\) is one of the things the
decomposition is designed to recover.

\end{orig}

\begin{rgt}
Lorem ipsum dolor sit amet, consectetur adipiscing elit, sed do eiusmod tempor incididunt ut labore et dolore magna aliqua.

\end{rgt}

\begin{orig}
Acute conditions tend toward TV-positive: a sprained ankle, a
post-operative pain syndrome, an episode of acute back pain all
resolve on a timescale of days to weeks even without
intervention, and a trial of an analgesic in such a setting will
see substantial \(TV\)-improvement in the placebo arm. Chronic
stable conditions (treated hypertension, controlled chronic
asthma, long-standing PTSD without active triggers) tend toward
\(TV\)-zero on the population mean but with substantial individual
variance. Chronic progressive conditions (Alzheimer's disease,
ALS, late-stage congestive heart failure, many cancers) tend
toward \(TV\)-negative: a trial of a neuroprotective agent in
early Alzheimer's must explicitly model the negative
natural-history trajectory, since the apparent treatment effect
is the slowing of progression, not absolute improvement.
Cyclical and triggered conditions (PTSD with trauma
anniversaries, seasonal-pattern major depression,
relapsing-remitting multiple sclerosis) show structured \(TV\)
that is neither monotone nor stationary, and may require
explicit covariates beyond the Gompertz form.

\end{orig}

\begin{rgt}
Lorem ipsum dolor sit amet, consectetur adipiscing elit, sed do eiusmod tempor incididunt ut labore et dolore magna aliqua.

\end{rgt}

\begin{orig}
Selection-induced \(TV\) is real and rarely zero, even when the
underlying condition is stable. Patients enrol at moments of
clinical engagement that often coincide with recent worsening, so
regression to the mean alone produces several points of apparent
\(TV\)-improvement in the early weeks of a trial. As the
introduction documents, the systematic-review evidence of\textsuperscript{\citeproc{ref-hrobjartsson2001}{11}},\textsuperscript{\citeproc{ref-hrobjartsson2004}{12}}, and\textsuperscript{\citeproc{ref-hrobjartsson2010}{13}} (the
last covering 202 trials across 60 conditions) and the
antidepressant analyses of\textsuperscript{\citeproc{ref-kirsch2008}{6}} and\textsuperscript{\citeproc{ref-locher2015}{7}} show that
much of what is conventionally attributed to placebo response is
indistinguishable from regression to the mean and natural-history
drift.

\end{orig}

\begin{rgt}
Lorem ipsum dolor sit amet, consectetur adipiscing elit, sed do eiusmod tempor incididunt ut labore et dolore magna aliqua.

\end{rgt}

\begin{orig}
Trial participation itself also induces \(TV\): the Hawthorne
effect, the structured-visit effect, and the diary-completion
effect appear in the data as \(TV\)-improvement that is
independent of any pharmacological or placebo response.

\end{orig}

\subsubsection{Can we assume no TV effect?}\label{can-we-assume-no-tv-effect}

\begin{rgt}
Lorem ipsum dolor sit amet, consectetur adipiscing elit, sed do eiusmod tempor incididunt ut labore et dolore magna aliqua.

\end{rgt}

\begin{orig}
But how safe is it to assume \(TV\) is zero? The short answer is:
rarely safe, and never safe without empirical support. The
longer answer is that \(TV\) may be assumed zero only if the
disease is in a stable chronic phase, the trial is short
relative to the disease timescale, the recruitment process does
not favor participants at peak severity, and the trial protocol
does not introduce ancillary care that itself improves symptoms.
For most chronic trials in psychiatric or pain conditions, at
least one of these conditions fails. The conservative default
is to estimate \(TV\) from the data, even at the cost of some
statistical efficiency, rather than assume it is zero.

\end{orig}

\begin{rgt}
Lorem ipsum dolor sit amet, consectetur adipiscing elit, sed do eiusmod tempor incididunt ut labore et dolore magna aliqua.

\end{rgt}

\begin{orig}
The decomposition makes this assumption testable. A model that
includes a \(TV\) component will return \(\hat{m}_{TV}\) with a
variance estimate; if both are small relative to \(\hat{m}_{BR}\)
and \(\hat{m}_{PB}\), the analyst has empirical support for the
simplifying assumption. If \(\hat{m}_{TV}\) is meaningfully large,
the assumption was wrong, and the trial needs the \(TV\) component
to produce unbiased inference about the others. Setting \(TV\) to
zero and absorbing it into \(\varepsilon\) is worse than estimating
it: the participant-level \(TV\) is structured (each participant
has their own slow trajectory) and not independent of the
on-drug-versus-off-drug schedule, so absorbing it into noise
produces residuals that are autocorrelated within participants
and heteroscedastic across them. The resulting standard errors
on the \(BR\) and \(PB\) estimates are wrong in both directions:
too small for participants whose \(TV\) is small, too large for
participants whose \(TV\) is large. Estimating \(TV\) explicitly
fixes both problems simultaneously.

\end{orig}

\subsubsection{Why a Gompertz, not a linear time covariate?}\label{why-a-gompertz-not-a-linear-time-covariate}

\begin{rgt}
Lorem ipsum dolor sit amet, consectetur adipiscing elit, sed do eiusmod tempor incididunt ut labore et dolore magna aliqua.

\end{rgt}

\begin{orig}
A natural objection is that a model with \texttt{week} as a
covariate already controls for time. But how adequate is such a
control? Three reasons suggest it is not enough. First, a
\texttt{week} covariate enforces a linear trajectory,
constraining each additional week to contribute a fixed amount
of natural-history change; real disease trajectories curve,
plateau, or accelerate, and the Gompertz form captures the
S-shaped or saturating patterns linear time cannot. Second, a
single \texttt{week} coefficient estimates a population-mean
trajectory that ignores the participant-level heterogeneity in
natural history. Third, \(BR\) and \(TV\) both evolve smoothly over
the same trial timescale, and a \texttt{week} covariate cannot
tell them apart from observational data alone; the design
contrasts discussed in section 4 are what supply
identifiability, and without the explicit decomposition the
\texttt{week} coefficient absorbs both.

\end{orig}

\section{Identifiability and trial-design contrasts}\label{identifiability-and-trial-design-contrasts}

\begin{rgt}
Lorem ipsum dolor sit amet, consectetur adipiscing elit, sed do eiusmod tempor incididunt ut labore et dolore magna aliqua.

\end{rgt}

\begin{orig}
The mathematical decomposition above is meaningless unless the
data contain enough information to pin down each component
separately. Trial design is what supplies that information. In
this section we shall characterize the hybrid
open-label-plus-blinded-discontinuation design, used as the
worked example throughout this paper, and show how each of its
three phases contributes a different combination of components
to the observed trajectory.

\end{orig}

{\def\LTcaptype{none} % do not increment counter
\begin{longtable}[]{@{}
  >{\raggedright\arraybackslash}p{(\linewidth - 6\tabcolsep) * \real{0.2500}}
  >{\raggedright\arraybackslash}p{(\linewidth - 6\tabcolsep) * \real{0.2500}}
  >{\raggedright\arraybackslash}p{(\linewidth - 6\tabcolsep) * \real{0.2500}}
  >{\raggedright\arraybackslash}p{(\linewidth - 6\tabcolsep) * \real{0.2500}}@{}}
\toprule\noalign{}
\begin{minipage}[b]{\linewidth}\raggedright
Phase
\end{minipage} & \begin{minipage}[b]{\linewidth}\raggedright
Drug?
\end{minipage} & \begin{minipage}[b]{\linewidth}\raggedright
Belief?
\end{minipage} & \begin{minipage}[b]{\linewidth}\raggedright
Components present
\end{minipage} \\
\midrule\noalign{}
\endhead
\bottomrule\noalign{}
\endlastfoot
Open-label on drug & Yes & Full & \(BR + PB(\eta = 1) + TV\) \\
Blinded discontinuation & Hidden & Hedged & \(BR \cdot \mathbf{1}_{\text{on drug}} + PB(\eta = 0.5) + TV\) \\
Open-label crossover & Either & Knows & \(BR \cdot \mathbf{1}_{\text{on drug}} + PB(\eta = 1) + TV\) \\
\end{longtable}
}

\begin{rgt}
Lorem ipsum dolor sit amet, consectetur adipiscing elit, sed do eiusmod tempor incididunt ut labore et dolore magna aliqua.

\end{rgt}

\begin{orig}
Three phase-by-allocation contrasts then identify the three
components. The open-label versus blinded contrast (within drug)
yields \(PB(\eta = 1) - PB(\eta = 0.5) \approx 0.5 \cdot PB\),
which isolates the belief component, but only conditional on a
known blinded-phase expectancy \(\eta\): from a two-phase design
just the product \(\eta \cdot m_{PB}\) is identified, and
separating \(\eta\) from \(m_{PB}\) requires the open-label placebo
phase. The blinded on-drug versus blinded placebo contrast
yields \(BR\) once drug carryover has elapsed, since \(PB\) and \(TV\)
match across allocation arms within the blinded phase. All
phases off drug yield \(TV\) plus any residual off-drug \(PB\), with
no active \(BR\) contribution beyond carryover and \(PB\) at the
off-drug expectation level. The \(\mathbf{1}_{\text{on drug}}\)
indicator in the phase table above is shorthand for the
continuous-decay drug state \(D_{it}\) defined in Table
\ref{tab:terms}.

This conditioning on \(\eta\) is the framework's principal
identification limitation and deserves emphasis rather than a
parenthetical. With a two-phase (open-label and blinded) design the
data identify only the product \(\eta \cdot m_{PB}\), so \(PB\) is
recovered up to the assumed blinded-phase expectancy; the working
value \(\eta \approx 0.5\) is a modeling assumption, not an estimate,
and absent a belief or expectancy measurement it is not falsifiable
from the outcome data alone. The identified set for \(m_{PB}\) is
accordingly the interval \(\{\widehat{\eta m_{PB}} / \eta : \eta \in
(0,1]\}\), which collapses to a point only when \(\eta\) is pinned down.
Two design features close the gap: an open-label placebo phase, which
supplies a second equation separating \(\eta\) from \(m_{PB}\), and direct
measurement of belief, as in the open-hidden and balanced-placebo
designs\textsuperscript{\citeproc{ref-colloca2004}{51},\citeproc{ref-rohsenow1981}{52}} that manipulate or record the
patient's treatment expectation. In a balanced-placebo design what
the patient is told and what the patient receives are crossed, so
that drug exposure and treatment belief vary independently. We therefore recommend that any trial
relying on the decomposition for a \(PB\)-versus-\(BR\) attribution either
include such a phase or collect an expectancy measure, and that
analyses report sensitivity of the attribution to \(\eta\) across its
plausible range.

\end{orig}

\begin{rgt}
Lorem ipsum dolor sit amet, consectetur adipiscing elit, sed do eiusmod tempor incididunt ut labore et dolore magna aliqua.

\end{rgt}

\begin{orig}
A trial that lacks any of these phases loses the corresponding
identifiability. A pure open-label design (no blinded
discontinuation) cannot separate \(BR\) from \(PB\), because both
are present at full strength throughout. A pure blinded design
(no open-label phase) cannot estimate \(\eta = 1\) \(PB\), because
no phase has belief at full strength. The decomposition
machinery is useful only to the extent that the design supplies
identifiability, and trial designers should evaluate any
proposed design by checking that each component has at least one
phase or contrast in which it is uniquely present.

\end{orig}

\begin{rgt}
Lorem ipsum dolor sit amet, consectetur adipiscing elit, sed do eiusmod tempor incididunt ut labore et dolore magna aliqua.

\end{rgt}

\begin{orig}
Identifiability in the design-contrast sense above is necessary
but not sufficient for stable estimation at finite sample sizes.
The three Gompertz trajectories share an onset shape and overlap
in calendar time, so their parameters can be collinear even when
the contrasts that identify them in principle are present. Two
conditions make the joint estimation credible: the design must
supply at least one phase in which each component varies while
the others are held fixed (the blinded-placebo contrast for
\(BR\), the open-label-versus-blinded contrast for \(PB\), and an
off-drug window for \(TV\)), and that off-drug window must be long
enough relative to the carryover half-life for \(TV\) to separate
from a still-decaying \(BR\). Where these conditions are weak the
components are jointly identified only up to a near-collinear
direction, and the symptom of that is the rank-deficiency the
pilot of section 7 encountered when the full phase-augmented
formula was fitted at \(N = 35\). We therefore do not assert that
the three components are always recoverable. The pre-registered
identifiability diagnostics of section 7 (convergence rates of
the component-matched fit, variance-inflation factors on the
component-specific Gompertz parameters, and the sensitivity of
the \(BR\) estimate to the assumed \(\eta(\text{phase})\) multipliers)
are the means by which the recoverable region of the
design-and-sample-size space is to be mapped, and the
framework's participant-level claims should be read as
conditional on falling inside that region.

\end{orig}

\section{Analysis-side methodology}\label{analysis-side-methodology}

\begin{rgt}
Lorem ipsum dolor sit amet, consectetur adipiscing elit, sed do eiusmod tempor incididunt ut labore et dolore magna aliqua.

\end{rgt}

\begin{orig}
In this section we shall take up the practical question of what
analysis model to use. A reasonable reaction to the formal model
in section 2 is the following: if we are simulating data with
three components plus noise, why is the analysis model written
as \texttt{Sx \textasciitilde\ bm + t + Dbc + bm:Dbc} with a
single random intercept and AR(1) residual correlation? Should
we not include three parametric Gompertz components in the
analysis, mirroring the DGP, so that the analyst sees the same
structure that generated the data? The question is reasonable.
Four constraints make a component-matched analysis impractical
for most aggregated N-of-1 trial sizes, and a modest extension
(a phase-by-treatment indicator) captures the practically
important part of the gap at minimal cost.

\end{orig}

\subsection{Four reasons not to match the DGP}\label{four-reasons-not-to-match-the-dgp}

\begin{rgt}
Lorem ipsum dolor sit amet, consectetur adipiscing elit, sed do eiusmod tempor incididunt ut labore et dolore magna aliqua.

\end{rgt}

\begin{orig}
The first reason is identifiability. The trial design must
supply the contrasts that identify each component, and many
designs supply only part of it. The decomposition is identified
by phase contrasts (open-label versus blinded for \(PB\), on-drug
versus blinded-placebo for \(BR\), off-drug timepoints for \(TV\)).
A pure open-label trial has none of these contrasts; in such a
trial \(BR\) and \(PB\) both rise from zero on treatment initiation
and are statistically indistinguishable from each other,
regardless of how the analyst writes the model. Forcing an
analysis to fit three components when the design supplies
contrasts for only one or two of them does not recover the
missing components; it produces unidentified or weakly
identified parameter estimates with implausibly tight standard
errors that misrepresent the analyst's actual inferential
precision.

\end{orig}

\begin{rgt}
Lorem ipsum dolor sit amet, consectetur adipiscing elit, sed do eiusmod tempor incididunt ut labore et dolore magna aliqua.

\end{rgt}

\begin{orig}
The second reason is degrees of freedom. Each component term
costs degrees of freedom that come out of the moderation
estimand. The biomarker-moderation question (the core scientific
aim of the predictive-biomarker validation program) is asked
through the \texttt{bm:Dbc} term. A component-matched analysis
adds a Gompertz \(BR\) (three parameters), a Gompertz \(PB\) (three
parameters plus a phase-specific expectancy multiplier), and a
Gompertz \(TV\) (three parameters), each with
participant-specific random effects on the maxima. That is on
the order of a dozen new parameters, half of them on a
non-linear scale. At trial sample sizes of \(N \in [30, 150]\),
the variance of \(\hat{\beta}_{bm:D}\) is dominated by participant
and timepoint counts, and adding parametric structure to absorb
residual confounding usually inflates the moderation estimand's
variance by more than it reduces bias.

\end{orig}

\begin{rgt}
Lorem ipsum dolor sit amet, consectetur adipiscing elit, sed do eiusmod tempor incididunt ut labore et dolore magna aliqua.

\end{rgt}

\begin{orig}
The third reason is convergence fragility. Component-matched
analyses are non-linear in the parameters and their convergence
is problematic in practice. Fitting Gompertz \(BR\), \(PB\), and
\(TV\) simultaneously requires a non-linear mixed model
(\texttt{nlme::nlme} with self-starting Gompertz forms), a
Bayesian hierarchical fit, or a structural-equation
approximation. Each is more sensitive to starting values, more
prone to local optima, and more likely to fail to converge on a
single trial dataset than the linear-mixed analogue. At
trial-relevant \(N\), the convergence rate of the
component-matched fit is typically below 90\%, and the
non-converging cells contaminate any aggregated analysis.

\end{orig}

\begin{rgt}
Lorem ipsum dolor sit amet, consectetur adipiscing elit, sed do eiusmod tempor incididunt ut labore et dolore magna aliqua.

\end{rgt}

\begin{orig}
The fourth reason is that the moderation estimand is
sufficiently robust to the BR-versus-PB split for the most
common scientific question. For predictive-biomarker validation,
the regulator wants to know whether the biomarker predicts
treatment response under the trial conditions, not a clean
pharmacological-versus-placebo decomposition of that prediction.
The simple \texttt{bm:Dbc} interaction is the right marginal
estimand for the validation question whether the moderation is
biological or expectation-mediated.

\end{orig}

\begin{rgt}
Lorem ipsum dolor sit amet, consectetur adipiscing elit, sed do eiusmod tempor incididunt ut labore et dolore magna aliqua.

\end{rgt}

\begin{orig}
It is worth stating the inferential target of the analysis in
the vocabulary of the estimands framework\textsuperscript{\citeproc{ref-iche9r1}{53}}, because
the biomarker-validation question is easy to pose ambiguously.
The primary estimand is the population-level
biomarker-treatment interaction, the \(bm{:}D_{bc}\)
coefficient: its population is the trial-eligible patient
population, its treatment contrast is on-drug versus off-drug
exposure summarized through the continuous indicator \(D_{bc}\),
its endpoint is the symptom trajectory, its population-level
summary is the interaction coefficient, and its intercurrent
events are handled by a treatment-policy strategy for dropout and
by the modeled exponential decay for residual carryover. The
pharmacological-component target \(\beta_{bm}^{BR}\) is a more
ambitious estimand: it is the biomarker-by-\(BR\) interaction that
would obtain with the belief-driven channel held at its off-drug
level, identified here through the phase contrasts of section 4
rather than by assumption. We flag explicitly that recovering a
participant-level interaction of this kind inherits the
difficulty Senn has emphasized for individual treatment effects\textsuperscript{\citeproc{ref-senn2004}{49},\citeproc{ref-senn2018}{54}}: a genuine participant-by-component signal
is separable from participant-by-period and within-person random
variation only when the design supplies repeated phase contrasts,
and the framework should be read as estimating these quantities
under that condition rather than as dissolving the well-known
problem of identifying individual effects from finite
within-person data.

\end{orig}

\subsection{A small extension that usually pays off: phase indicators}\label{a-small-extension-that-usually-pays-off-phase-indicators}

\begin{rgt}
Lorem ipsum dolor sit amet, consectetur adipiscing elit, sed do eiusmod tempor incididunt ut labore et dolore magna aliqua.

\end{rgt}

\begin{orig}
The smallest useful step toward component-aware analysis,
without paying the costs above, is to add a phase indicator
and its interaction with the drug indicator. Schematically,

\begin{verbatim}
Sx ~ bm + t + Dbc + phase + Dbc:phase
       + bm:Dbc + bm:Dbc:phase
       + (1 | ptID),  correlation = corCAR1(...)
\end{verbatim}

Here \texttt{phase} is a categorical covariate (open-label,
blinded, post-crossover) that absorbs the systematic differences
in \(PB\) expectancy across phases. \texttt{Dbc:phase} tests
whether the apparent drug effect shrinks under blinding, a
signature that the open-label drug effect was partly
\(PB\)-mediated. \texttt{bm:Dbc:phase} tests whether the
moderation itself shrinks under blinding, a signature that the
biomarker-treatment interaction was partly
biomarker-by-\(PB\) rather than biomarker-by-\(BR\).

\end{orig}

\begin{rgt}
Lorem ipsum dolor sit amet, consectetur adipiscing elit, sed do eiusmod tempor incididunt ut labore et dolore magna aliqua.

\end{rgt}

\begin{orig}
This extension adds three to four parameters at most, well
within budget at any realistic \(N\). It does not separately
estimate \(BR\) and \(PB\) as components, but it does identify the
change in apparent drug effect that occurs when belief is
hedged. For most analyses, this is the practically relevant
question, and it can be answered without committing to the
Gompertz parametric form. Where both the one-component and the
phase-augmented analyses are reported, the additional
\texttt{Dbc:phase} and \texttt{bm:Dbc:phase} terms introduce a
small family of belief-attenuation tests, and the attribution
inference they support should be treated as a pre-specified
secondary analysis with its multiplicity acknowledged rather than
as a second bite at the primary interaction test.

\end{orig}

\subsection{When to fit the full decomposition}\label{when-to-fit-the-full-decomposition}

\begin{rgt}
Lorem ipsum dolor sit amet, consectetur adipiscing elit, sed do eiusmod tempor incididunt ut labore et dolore magna aliqua.

\end{rgt}

\begin{orig}
A component-matched analysis merits the effort when sample size
is in the several hundreds (a phase-3 trial fielded as
aggregated N-of-1, or a meta-analysis pooling multiple smaller
trials); the trial design provides timepoints in every phase
combination, with enough density per cell to constrain the
Gompertz curves; the scientific question requires the BR-PB
split (a pharmacology paper claiming a specific mechanism, where
ruling out a placebo contribution is the central point); and a
Bayesian hierarchical fit with informative priors on the
Gompertz parameters is feasible. In that regime, \texttt{lcmm},
\texttt{nlme::nlme} with self-starting Gompertz forms, or a
\texttt{brms} hierarchical formulation can return
participant-specific \(BR\), \(PB\), and \(TV\) component estimates
that mirror the DGP. Outside that regime, the linear-mixed
analysis with phase indicators is the better tool.

\end{orig}

\begin{rgt}
Lorem ipsum dolor sit amet, consectetur adipiscing elit, sed do eiusmod tempor incididunt ut labore et dolore magna aliqua.

\end{rgt}

\begin{orig}
The asymmetry between DGP and analysis is not a flaw of the
framework; it is a deliberate consequence of the asymmetry
between simulation (where the analyst controls the DGP and can
specify any structure, however elaborate) and inference (where
the analyst must fit a model to a finite dataset and pay the
identifiability cost of every additional parameter). Simulation
studies use the rich DGP to characterize how a simpler analysis
behaves under realistic data generation; they do not require the
analysis to match the DGP in order to produce useful inference.

\end{orig}

\begin{rgt}
Lorem ipsum dolor sit amet, consectetur adipiscing elit, sed do eiusmod tempor incididunt ut labore et dolore magna aliqua.

\end{rgt}

\begin{orig}
A word on the carryover parameterization is in order, since it
is the one place where the analysis model makes a strong
functional-form commitment. The continuous drug indicator
\(D_{bc}\) decays as \(\exp(-\lambda t_{sd})\) off drug, a
one-parameter exponential governed by the carryover half-life.
This is deliberately parsimonious and is best understood as a
constrained special case of the more flexible carryover models
now available for repeated-measures and N-of-1 data. Liao et al.\textsuperscript{\citeproc{ref-liao2023}{31}} model carryover through a Bayesian distributed-lag
structure with autocorrelated errors, relaxing the single-rate
exponential assumption at the cost of additional parameters;
Daza\textsuperscript{\citeproc{ref-dazaCausalAnalysisSelftracked2018}{55}} formalizes an average-period treatment-effect
estimand built to accommodate autocorrelation, trend, and slow
carryover jointly; and G\{"a\}rtner et al.\textsuperscript{\citeproc{ref-gaertnerCarryover2023}{56}}
benchmark linear models against G-estimation and Bayesian-network
alternatives specifically under carryover. We adopt the
exponential form because it ties directly to a pharmacokinetic
half-life and conserves degrees of freedom for the moderation
estimand at trial-relevant sample sizes, but the choice is an
assumption rather than a finding. The pre-registered simulation
program of section 7 should accordingly include a
distributed-lag analysis arm, so that the cost of the
exponential-decay restriction is quantified against the richer
alternative rather than asserted.

\end{orig}

\section{Worked example: information loss without decomposition}\label{worked-example-information-loss-without-decomposition}

\subsection{What the one-component analysis estimates}\label{what-the-one-component-analysis-estimates}

\begin{rgt}
Lorem ipsum dolor sit amet, consectetur adipiscing elit, sed do eiusmod tempor incididunt ut labore et dolore magna aliqua.

\end{rgt}

\begin{orig}
Before the numerical illustration it is worth establishing
algebraically what the one-component analysis actually estimates,
because the qualitative failure modes of section 1 are
consequences of an identity rather than empirical tendencies that
require simulation to demonstrate. Write the within-participant
symptom reduction as
\(R_{it} = \mathrm{BL}_i - Y_{it} = BR_{it} + PB_{it} + TV_{it} -
\varepsilon_{it}\). The one-component analysis fits a single drug
indicator \(D_{it}\),

\[
R_{it} \;=\; \alpha + \beta_D\,D_{it} + g(t) + u_i + e_{it},
\]

and reports \(\hat\beta_D\) as `the treatment effect'. Taking the
on-drug minus off-drug contrast that \(\beta_D\) targets, and using
that the natural-history component is by construction independent
of treatment status so that its on/off difference cancels in a
within-participant comparison, the probability limit of the
estimator is

\[
\operatorname{plim}\hat\beta_D \;=\; m_{BR} + b_{PB} + b_{TV},
\]

where \(b_{PB} = E[PB \mid \text{on}] - E[PB \mid \text{off}]\) is
the belief differential between the contrasted phases and \(b_{TV}\)
collects any natural-history contribution that fails to cancel
under the contrast actually used; \(b_{TV}\) is zero for a clean
on/off within-participant contrast and strictly positive for a
baseline-anchored pre-post analysis, in which case
\(b_{PB} = E[PB]\) and \(b_{TV} = E[TV]\). To fix which case is in
play below: the one-component analysis used in the worked example
and the simulations of this paper is the design-contrast
(drug-on/off) form, so \(b_{TV} \approx 0\) there and the operative
displacement is the belief term \(b_{PB}\); the baseline-anchored
pre-post form, which additionally carries \(b_{TV} = E[TV]\), is the
worse case retained only for comparison. Both terms are generally
non-negative for an improving outcome, so the one-component estimate
over-states the pharmacological component \(m_{BR}\). The
displacement is a property of the probability limit and is
unchanged as \(N \to \infty\): a larger one-component trial
converges on the same biased value. This is the precise sense in
which the placebo-attribution failure of section 1 is a problem
of identification, not of precision.

\end{orig}

\begin{rgt}
Lorem ipsum dolor sit amet, consectetur adipiscing elit, sed do eiusmod tempor incididunt ut labore et dolore magna aliqua.

\end{rgt}

\begin{orig}
The same algebra settles the biomarker-validation failure. The
precision-medicine analysis regresses a participant-level response
summary on a baseline biomarker \(B_i\). If that summary is the
one-component per-participant effect, which estimates
\(m_{BR,i} + w_{PB}\,m_{PB,i} + w_{TV}\,m_{TV,i}\) for
design-determined loadings \(w_{PB}, w_{TV} \ge 0\), then by the
linearity of covariance the fitted biomarker slope has probability
limit

\[
\beta_{bm}^{\text{lumped}}
\;=\; \frac{\operatorname{Cov}(m_{BR} + w_{PB} m_{PB} +
  w_{TV} m_{TV},\, B)}{\operatorname{Var}(B)}
\;=\; \beta_{bm}^{BR} + w_{PB}\,\beta_{bm}^{PB} +
  w_{TV}\,\beta_{bm}^{TV},
\]

where \(\beta_{bm}^{C} = \operatorname{Cov}(m_C, B) /
\operatorname{Var}(B)\) is the component-specific slope. A biomarker
correlated with placebo responsiveness or with the natural-history
trajectory, so that \(\beta_{bm}^{PB}\) or \(\beta_{bm}^{TV}\) is
non-zero, therefore registers a non-zero one-component slope even when its
pharmacological association \(\beta_{bm}^{BR}\) is exactly zero. This
is the biomarker-validation failure of section 1, and it too is an
identity under the additive decomposition rather than an empirical
tendency: no sample size separates the one-component slope from a genuine
pharmacological one, because the two coincide in the estimand
whenever \(w_{PB}\beta_{bm}^{PB} + w_{TV}\beta_{bm}^{TV} \ne 0\).

We make two qualifications explicit, because they bound how much
weight the framework can carry. First, both displays are immediate
corollaries of the additive decomposition and the linearity of the
projection that defines the estimator; they are not a derived result
in any deeper sense, and their value is in naming the displacement
terms, not in the algebra. Second, they assume the components combine
additively. Under the subadditive interaction documented in section
3.3, where the joint effect of \(BR\) and \(PB\) falls below their sum, a
first-order correction proportional to \(\operatorname{Cov}(m_{BR}\cdot
m_{PB}, B)\) enters the biomarker slope, and the identities above
should be read as the leading behavior in the additive case rather
than as exact under interaction. The empirical magnitude of that
correction term is itself an open question that the framework's
\(\gamma\)-interaction simulation (section 3.3) is designed to address.

\end{orig}

\begin{rgt}
Lorem ipsum dolor sit amet, consectetur adipiscing elit, sed do eiusmod tempor incididunt ut labore et dolore magna aliqua.

\end{rgt}

\begin{orig}
Two consequences for how the framework should be evaluated follow.
First, the quantity that carries the bias is the location of the
estimand, \(\operatorname{plim}\hat\beta_D\), not its sampling
variability: the one-component coefficient remains a well-estimated number
with a valid Wald test of its own nullity, and that test rejects at
the rate dictated by the displaced value. A significance test of
the one-component coefficient therefore cannot, in principle, detect the
bias, because the bias does not move the coefficient toward zero.
Detecting it requires comparing \(\hat\beta_D\) against \(m_{BR}\),
which is exactly the comparison the decomposition, or a phase
contrast that isolates \(m_{BR}\), makes available and the one-component
analysis does not. Second, and as a corollary for the simulation
evidence, a rejection-rate or power summary is the wrong instrument
for this bias: the right target is the coefficient itself and its
offset from \(m_{BR}\). This is why the pilot of section 7.1 found
the biomarker-treatment rejection rate essentially flat across
the placebo-strength axis even though that axis is precisely what
drives the predicted bias; the prediction lives in the mean of
\(\hat\beta_{bm:D_{bc}}\), and a production run that means to test it
must read the bias off the coefficient, not off the test.

\end{orig}

\subsection{A numerical illustration}\label{a-numerical-illustration}

\begin{rgt}
Lorem ipsum dolor sit amet, consectetur adipiscing elit, sed do eiusmod tempor incididunt ut labore et dolore magna aliqua.

\end{rgt}

\begin{orig}
To illustrate the inferential cost of failing to decompose, we
shall work through a simple numerical example. Consider two
participants, A and B, enrolled in a sixteen-week N-of-1 trial
of prazosin for PTSD. Both are otherwise similar: same age,
same baseline severity, same trial design. They differ in two
latent traits that the trial does not measure directly.
Participant A is a high placebo responder and a moderate
pharmacological responder, with true component values
\(m_{BR}^A = 4\), \(m_{PB}^A = 6\), \(m_{TV}^A = 1\). Participant B
is a low placebo responder and a strong pharmacological
responder, with \(m_{BR}^B = 8\), \(m_{PB}^B = 1\), \(m_{TV}^B = 1\).

\end{orig}

\begin{rgt}
Lorem ipsum dolor sit amet, consectetur adipiscing elit, sed do eiusmod tempor incididunt ut labore et dolore magna aliqua.

\end{rgt}

\begin{orig}
In the open-label phase (weeks 1 to 8), both \(BR\) and \(PB\) run
at full strength and \(TV\) accumulates. By week 8 both
participants have reached close to their saturation values:

{\def\LTcaptype{none} % do not increment counter
\begin{longtable}[]{@{}lllll@{}}
\toprule\noalign{}
& \(BR\) & \(PB(\eta = 1)\) & \(TV\) & Total improvement \\
\midrule\noalign{}
\endhead
\bottomrule\noalign{}
\endlastfoot
Participant A & 4.0 & 6.0 & 1.0 & \textbf{11.0} \\
Participant B & 8.0 & 1.0 & 1.0 & \textbf{10.0} \\
\end{longtable}
}

A clinician examining only the totals would conclude that the
two participants are responding similarly: both improved by
about ten points, both look like clinical successes. A simple
t-test on the open-label change would estimate a population
drug effect of about ten points and would not distinguish the
two participants.

\end{orig}

\begin{rgt}
Lorem ipsum dolor sit amet, consectetur adipiscing elit, sed do eiusmod tempor incididunt ut labore et dolore magna aliqua.

\end{rgt}

\begin{orig}
In the blinded discontinuation phase (weeks 9 to 12), both
participants are silently switched to placebo at the start of
week 9. The expectancy multiplier drops from 1.0 to 0.5. For
clarity of exposition we take the carryover half-life to be short
relative to the four-week blinded window, so that \(BR\) has
effectively cleared by the end of week 12; the more realistic
case of a slowly decaying \(BR\), modeled by the carryover
machinery of section 3.1, attenuates but does not eliminate the
contrast below. By the end of week 12, the components are:

{\def\LTcaptype{none} % do not increment counter
\begin{longtable}[]{@{}lllll@{}}
\toprule\noalign{}
& \(BR\) & \(PB(\eta = 0.5)\) & \(TV\) & Total improvement \\
\midrule\noalign{}
\endhead
\bottomrule\noalign{}
\endlastfoot
Participant A & 0.0 & 3.0 & 1.0 & \textbf{4.0} \\
Participant B & 0.0 & 0.5 & 1.0 & \textbf{1.5} \\
\end{longtable}
}

The two participants now look quite different. Participant A
still shows a sizeable improvement (four points) under blinded
placebo; participant B has nearly returned to baseline (one and
a half points). This is the diagnostic contrast that the
decomposition exploits. For simplicity the table holds \(TV\) at
its week-8 value; in reality \(TV\) continues to accumulate over
weeks 9 to 12, which shifts both totals by a common amount and
leaves the \(BR\)-versus-\(PB\) contrast on which the example turns
unchanged.

\end{orig}

\begin{rgt}
Lorem ipsum dolor sit amet, consectetur adipiscing elit, sed do eiusmod tempor incididunt ut labore et dolore magna aliqua.

\end{rgt}

\begin{orig}
A one-component fixed-effects regression of total symptom
change on phase-by-allocation indicators sees:

\begin{itemize}
\tightlist
\item
  Open-label means: \(11.0\) (A), \(10.0\) (B).
\item
  Blinded-placebo means: \(4.0\) (A), \(1.5\) (B).
\item
  Estimated drug effect: \(11.0 - 4.0 = 7.0\) (A), \(10.0 - 1.5 =
  8.5\) (B). Population mean: \(7.75\).
\item
  Estimated placebo effect: not identifiable from this contrast.
\item
  Estimated natural history: not identifiable from this contrast.
\end{itemize}

The one-component model has produced a single number conflated
with the placebo response that washes out under blinding. It
cannot say whether the drug effect is the same for participants
A and B; it estimates `drug' at 7.75 points on average, which
is an average of the true 4.0 (A) and 8.0 (B) mixed with the
placebo wash-out, with no principled way to decompose further.

\end{orig}

\begin{rgt}
Lorem ipsum dolor sit amet, consectetur adipiscing elit, sed do eiusmod tempor incididunt ut labore et dolore magna aliqua.

\end{rgt}

\begin{orig}
A linear-mixed-effects analysis with the full BR-PB-TV
decomposition, fitted to the same data, can recover the
individual components by exploiting the differential expectancy
between phases and the on-drug-versus-blinded-placebo contrast.
After fitting, the model returns \(\hat{m}_{BR}^A = 4.0\),
\(\hat{m}_{BR}^B = 8.0\), \(\hat{m}_{PB}^A = 6.0\),
\(\hat{m}_{PB}^B = 1.0\), \(\hat{m}_{TV}^A = \hat{m}_{TV}^B = 1.0\)
(at the population point estimate; participant-level estimates
have meaningful uncertainty at \(N\) in the trial-relevant range,
which the simulation study below characterizes).

\end{orig}

\begin{rgt}
Lorem ipsum dolor sit amet, consectetur adipiscing elit, sed do eiusmod tempor incididunt ut labore et dolore magna aliqua.

\end{rgt}

\begin{orig}
The information loss from combining the components into one indicator is not abstract.
We may compare the inferences directly:

{\def\LTcaptype{none} % do not increment counter
\begin{longtable}[]{@{}
  >{\raggedright\arraybackslash}p{(\linewidth - 4\tabcolsep) * \real{0.3333}}
  >{\raggedright\arraybackslash}p{(\linewidth - 4\tabcolsep) * \real{0.3333}}
  >{\raggedright\arraybackslash}p{(\linewidth - 4\tabcolsep) * \real{0.3333}}@{}}
\toprule\noalign{}
\begin{minipage}[b]{\linewidth}\raggedright
Question
\end{minipage} & \begin{minipage}[b]{\linewidth}\raggedright
One-component answer
\end{minipage} & \begin{minipage}[b]{\linewidth}\raggedright
Three-component answer
\end{minipage} \\
\midrule\noalign{}
\endhead
\bottomrule\noalign{}
\endlastfoot
Does the drug work for A? & \textasciitilde7.75 points (one-component). & \(BR_A = 4.0\), distinct from her larger placebo response. \\
Does the drug work for B? & \textasciitilde7.75 points (one-component). & \(BR_B = 8.0\), twice the magnitude of A. \\
Are A and B different responders? & No discernible difference. & B has 2x A's \(BR\) but 1/6 of A's \(PB\). \\
Is a biomarker predicting BR useful here? & Cannot answer. & Yes if the biomarker correlates with \(\hat{m}_{BR}\). \\
Should A be continued long-term? & Cannot say. & Probably yes (\(BR = 4\) is real and pharmacological), but the larger placebo component will not persist if expectations decay. \\
Population mean drug effect? & 7.75 points (biased). & \(6.0 = (4.0 + 8.0)/2\) (correct mean of \(BR\)). \\
\end{longtable}
}

The one-component model has produced one biased average where
the three-component model has produced six clinically actionable
quantities. Predictive-biomarker validation, which is a central
scientific aim of the trial design family considered here,
becomes impossible without the decomposition.

\end{orig}

\begin{rgt}
Lorem ipsum dolor sit amet, consectetur adipiscing elit, sed do eiusmod tempor incididunt ut labore et dolore magna aliqua.

\end{rgt}

\begin{orig}
A second example at the population level illustrates the same
problem. Consider two cohorts of N-of-1 participants enrolled
in methodologically identical trials of the same drug. Cohort 1
(early trial, naive patients) has mean \(m_{BR} = 5\), mean
\(m_{PB} = 7\), mean \(m_{TV} = 1\). Cohort 2 (later trial,
patients who have had multiple prior failed treatments and are
skeptical) has mean \(m_{BR} = 5\), mean \(m_{PB} = 2\), mean
\(m_{TV} = 1\). By construction the drug works equally well in
both cohorts. But the open-label totals differ substantially:
cohort 1 reports thirteen-point improvements on average, cohort
2 reports eight-point improvements. A one-component analysis
comparing the two cohorts would conclude that the drug works
less well in cohort 2 and might trigger a search for
confounders or sub-population effects. None of this is real;
the difference is entirely in \(PB\), specifically in the
population baseline level of placebo responsiveness. The
three-component decomposition, applied to either cohort, returns
\(\hat{m}_{BR} = 5\) in both. Unfortunately, this is the kind of
apparent non-replicability that pharmacology has spent the last
decade struggling to navigate; the decomposition is one of the
formal tools that makes it tractable.

\end{orig}

\section{Simulation study}\label{simulation-study}

\begin{rgt}
Lorem ipsum dolor sit amet, consectetur adipiscing elit, sed do eiusmod tempor incididunt ut labore et dolore magna aliqua.

\end{rgt}

\begin{orig}
The simulation study reported and pre-registered in this section
is structured under the\textsuperscript{\citeproc{ref-morris2019simulation}{57}} ADEMP framework, specifying
Aims, Data-generating mechanisms, Estimands, Methods, and
Performance measures. The pre-registered plan is at
\texttt{analysis/scripts/component-decomposition/00-ademp-pre-registration.md}.
A Monte Carlo simulation calibrated to the prazosin/PTSD
reference parameters quantifies the bias and identifiability
properties of the three analysis strategies described in
section 5 across a parameter grid.

\end{orig}

\subsection{Pilot validation of analysis machinery (100 replicates, 6 cells)}\label{pilot-validation-of-analysis-machinery-100-replicates-6-cells}

\begin{rgt}
Lorem ipsum dolor sit amet, consectetur adipiscing elit, sed do eiusmod tempor incididunt ut labore et dolore magna aliqua.

\end{rgt}

\begin{orig}
The production simulation program is the empirical core of this
manuscript; its first realized run, Study A, is reported in section
7.3, while the fuller pre-registered grid (section 7.2) remains
partly outstanding. The pilot reported here precedes both and
confirms only that the analysis machinery runs end-to-end.

\end{orig}

\begin{rgt}
Lorem ipsum dolor sit amet, consectetur adipiscing elit, sed do eiusmod tempor incididunt ut labore et dolore magna aliqua.

\end{rgt}

\begin{orig}
A 100-replicate-per-cell pilot was run at the reference Hybrid
open-label-plus-blinded-discontinuation (OL+BDC) design, \(N = 35\), on the 6-cell slice
\(pb\_strength \in \{0, 6.5, 13\} \times\) analysis \(\in\)
\{one-component, phase-augmented\}, using 8-core
\texttt{parallel::mclapply}. The convergence rate was 1.0
across all six cells. Per-replicate output is checkpointed at
\texttt{analysis/data/quick-sim/06-decomp-replicates.rds} and
the Morris-format cell-level summary at
\texttt{analysis/data/quick-sim/06-decomp-summary.txt}. The
exploratory summaries are diagnostic only. The one-component
\texttt{bm:Dbc} test showed power \(0.95\), \(0.93\), and \(0.83\)
across \(pb\_strength \in \{0, 6.5, 13\}\), while the
phase-augmented analysis showed \(0.65\), \(0.65\), and \(0.61\); the
mean estimated \(\hat\beta_{bm:D_{bc}}\) was essentially flat
within each analysis (about \(-2.25\) for the one-component and
\(-2.21\) for the phase-augmented fit). At 100 replicates the
MCSE on a power estimate near \(0.7\) is
about \(0.046\), so these figures cannot by themselves separate
the two analyses; their interpretation, including why the
phase-augmented analysis loses its bias-detection contrast at
this \(N\), is resolved by the production run of section 7.3. No confirmatory bias or
power claim is drawn from the pilot, and all substantive claims
are deferred to the production program below.

\end{orig}

\begin{rgt}
Lorem ipsum dolor sit amet, consectetur adipiscing elit, sed do eiusmod tempor incididunt ut labore et dolore magna aliqua.

\end{rgt}

\begin{orig}
Two design-level observations from the pilot inform the
production specification. First, at \(N = 35\) the rank-deficiency
of the full phase-augmented formula required dropping the
\(D_{bc}{:}\text{phase}\) interaction term, removing the precise
contrast that would have isolated the
biomarker-treatment-by-phase signal the framework is designed
to detect. Second, the predicted bias of the one-component
analysis arises in the regression coefficient
\(\hat\beta_{bm:D_{bc}}\), not in the indicator-level rejection
rate, so power on the biomarker-treatment Wald test is not
the right summary statistic to detect the framework's predicted
bias. The production program below addresses both observations
through larger \(N\), a richer trial design that admits the full
phase-augmented formula, and reporting of the coefficient mean
against the implied true value rather than the rejection rate
alone.

\end{orig}

\subsection{Production-simulation specification}\label{production-simulation-specification}

\begin{rgt}
Lorem ipsum dolor sit amet, consectetur adipiscing elit, sed do eiusmod tempor incididunt ut labore et dolore magna aliqua.

\end{rgt}

\begin{orig}
The production program follows the\textsuperscript{\citeproc{ref-morris2019simulation}{57}} ADEMP framework and
targets four deliverables: the bias of \(\hat\beta_{bm:D_{bc}}\) for the
one-component, phase-augmented, and full-decomposition analyses across
\(N \in \{35,70,200,300\}\) and a \(pb\_strength\) grid; identifiability
metrics for the full-decomposition fit (convergence, variance-inflation
factors, and sensitivity of the \(BR\) estimate to the
\(\eta(\text{phase})\) multipliers); Type I error and power for the
interaction test under an \(\eta\)-mismatch, \(TV\)-magnitude, and dropout
sensitivity grid, with null cells at \(n_{\text{reps}} = 5{,}000\) for
MCSE \(0.003\); and the participant-level interaction signal recovered as
a function of \((m_{BR}, m_{PB}, m_{TV})\), which the one-component-total
regression does not preserve. The \texttt{tidyverse} implementation
collection drives the grid; driver scripts, checkpointed cell-level
summaries, per-replicate output, and the committed ADEMP
pre-registration all reside under
\texttt{analysis/scripts/component-decomposition/}. The realized Study A
run reported next used 1,000 replicates per alternative cell.

Two features of the sample-size axis should be read together with the
results. First, \(N\) is the total across randomization paths, following
the compendium convention, but the driver fixes the count passed to
each path: the two smaller cells use a single-path design, so their
totals are 35 and 70 as the driver states them, while the two larger
cells use a two-path design and therefore analyze 200 and 300
participants where the driver grid records 100 and 150. The totals are
quoted here. Second, that same threshold switches the design itself,
from the eight-visit single-path form to the sixteen-visit two-path
form with the crossover phase. The axis is consequently not a clean
sample-size sweep, and the step between the 70 and 200 cells combines
a larger sample, more visits per participant, and an added crossover
contrast. Comparisons within either pair are interpretable; a
comparison across the threshold is not attributable to sample size
alone.

\end{orig}

\subsection{Study A results: bias and power at 1,000 replicates per alternative cell}\label{study-a-results-bias-and-power-at-1000-replicates-per-alternative-cell}

\begin{rgt}
Lorem ipsum dolor sit amet, consectetur adipiscing elit, sed do eiusmod tempor incididunt ut labore et dolore magna aliqua.

\end{rgt}

\begin{orig}
The production Study A run (1,000 replicates per alternative cell,
5,000 per null cell, 80 alternative cells crossing
\(m_{PB} \in \{0,1,3,6,10\}\), \(m_{TV} \in \{-1,0,1,2\}\), and
\(N \in \{35,70,200,300\}\), plus four null cells) returns a result
that does not support the framework's central bias prediction.
The DGP couples the baseline biomarker to the
\(BR\) component only, through the on-drug biomarker-\(BR\) covariance
channel, so the true biomarker-treatment effect is the
biomarker-by-\(BR\) slope,
\(\beta_{bm:D_{bc}} = -c_{bm}\,\sigma_{BR}/\sigma_{bm} = -2.25\), and
the biomarker is by construction uncorrelated with \(PB\) and \(TV\).
Under the one-component analysis the mean estimated
\(\hat\beta_{bm:D_{bc}}\) tracks this value across the whole grid:
the mean absolute bias is \(0.016\) at \(N=35\), \(0.016\) at \(N=70\),
\(0.006\) at \(N=200\), and \(0.006\) at \(N=300\), in every case within
one MCSE (respectively about \(0.028\),
\(0.019\), \(0.009\), and \(0.007\)) of zero, with no monotone trend
across \(m_{PB}\) or \(m_{TV}\). The predicted placebo-attribution
bias of the biomarker-treatment estimate does not appear.

This is not a failure of the simulation; it is the behavior the
identity of section 6.1 requires. That identity writes the one-component
biomarker slope as \(\beta_{bm}^{BR} + w_{PB}\beta_{bm}^{PB} +
w_{TV}\beta_{bm}^{TV}\). With the biomarker coupled to \(BR\) alone,
\(\beta_{bm}^{PB} = \beta_{bm}^{TV} = 0\), so the one-component slope equals
the pharmacological slope no matter how large \(m_{PB}\) or \(m_{TV}\)
grows. Study A varies the magnitudes of the unmodeled components,
which the identity shows to be the innocuous axis; the axis that
produces biomarker-validation bias is the biomarker's correlation
with \(PB\) or \(TV\), which this DGP holds at zero. The clean null is
therefore a confirmation of the identity and a correction to the
looser claim of section 1 that large \(PB\) or \(TV\) by itself biases
biomarker validation. It does not; biomarker-component coupling
does.

\end{orig}

\begin{rgt}
Lorem ipsum dolor sit amet, consectetur adipiscing elit, sed do eiusmod tempor incididunt ut labore et dolore magna aliqua.

\end{rgt}

\begin{orig}
The comparison between analyses is equally clear. The
phase-augmented analysis is not less biased than the one-component
analysis at any sample size (its mean absolute bias is \(0.028\)
versus \(0.016\) at \(N=35\) and is comparable elsewhere), and it is
markedly less powerful for the biomarker-treatment test: mean
power across the grid is \(0.35\) versus \(0.59\) at \(N=35\) and \(0.62\)
versus \(0.90\) at \(N=70\), with both analyses reaching essentially
full power by \(N=200\). Because the DGP contains no biomarker-\(PB\)
contamination for the phase contrast to remove, the extra
phase-by-treatment terms add estimation cost without buying any
bias reduction, and the analysis is dominated. Convergence was
approximately complete for both analyses at every sample size
(minimum convergence rate \(0.998\)), so the rank-deficiency that
forced dropping the \texttt{Dbc:phase} term in the \(N=35\) pilot,
which we had advanced as the explanation for the pilot's inverted
power ranking, is not operative in the production run: the inverted
ranking persists with the full formula fitted and converged. The
pilot's power inversion was therefore a genuine feature of this
DGP, not an artifact of rank-deficiency.

Calibration is adequate for both analyses. Nominal-95\%
confidence-interval coverage of \(\beta_{bm:D_{bc}}\) ranges from
\(0.96\) to \(0.99\) for the one-component analysis (slightly
conservative) and from \(0.94\) to \(0.98\) for the phase-augmented
analysis. Type I error at the four null cells is controlled,
ranging from \(0.031\) to \(0.056\) for the one-component analysis and
\(0.038\) to \(0.049\) for the phase-augmented analysis; the
one-component estimate at \(N=200\) (\(0.056\), MCSE \(0.003\)) sits
1.9 standard errors above nominal but is within the simulation's
Monte Carlo uncertainty.

\end{orig}

\begin{rgt}
Lorem ipsum dolor sit amet, consectetur adipiscing elit, sed do eiusmod tempor incididunt ut labore et dolore magna aliqua.

\end{rgt}

\begin{orig}
Two conclusions follow for the Study A program. First, the
result sharpens rather than supports the section 1 framing: the
biomarker-validation failure is real as an identity, but it is
driven by the biomarker's coupling to \(PB\) or \(TV\), not by the
magnitude of those components, and Study A as pre-registered
varies the latter while holding the former at zero. The decisive
axis is the biomarker-\(PB\) correlation \(c_{bm,PB}\), which the
contaminated-biomarker pilot in the following subsection varies
directly. Second, on the present evidence the phase-augmented
analysis carries a real power cost without a compensating bias
reduction under an uncontaminated biomarker, so it should not be
recommended as a default for this design; the pilot below
evaluates whether contamination changes this verdict.

\end{orig}

\subsection{Contaminated-biomarker pilot (100 replicates per cell)}\label{contaminated-biomarker-pilot-100-replicates-per-cell}

\begin{rgt}
Lorem ipsum dolor sit amet, consectetur adipiscing elit, sed do eiusmod tempor incididunt ut labore et dolore magna aliqua.

\end{rgt}

\begin{orig}
A supplementary pilot varying the biomarker-\(PB\) correlation
\(c_{bm,PB} \in \{0, 0.15, 0.30, 0.45\}\) at two \(PB\) magnitude
levels (\(m_{PB} \in \{0, 6\}\)) and two sample sizes
(\(N \in \{70, 300\}\)), using 100 replicates per cell, confirms
the identity-based prediction and reveals an unexpected structural
result.

At \(N=300\) the one-component bias scales monotonically with
\(c_{bm,PB}\): mean \(\hat\beta_{bm:D_{bc}}\) moves from \(-2.31\)
(\(c_{bm,PB}=0\), bias \(-0.06\)) to \(-2.37\) (\(c_{bm,PB}=0.15\),
bias \(-0.12\)) to \(-2.08\) (\(c_{bm,PB}=0.30\), bias \(+0.17\)) to
\(-1.74\) (\(c_{bm,PB}=0.45\), bias \(+0.51\)). The last value sits
30 MCSEs from the true \(-2.25\), establishing
the effect unambiguously. The direction of bias is positive
(toward zero): high-\(c_{bm,PB}\) participants experience stronger
\(PB\), which in the open-label phase adds to the apparent drug
effect, but this elevation is partially absorbed by the \(bm\) main
effect, leaving the \(bm:D_{bc}\) interaction coefficient attenuated
relative to the pure pharmacological value.

The unexpected result is that \(m_{PB}\) does not modulate the
bias: at the contaminated cells (\(c_{bm,PB} \geq 0.30\), \(N=300\))
the \(m_{PB}=0\) and \(m_{PB}=6\) biases are identical to within
\(0.006\), well inside Monte Carlo error, even as the bias itself
grows to \(+0.51\); the only larger discrepancy is at
\(c_{bm,PB}=0\), where both biases are near zero and differ by
\(0.067\) at the pilot's 100-replicate resolution. The driver is
therefore \(c_{bm,PB} \times \sigma_{PB}\), not
\(c_{bm,PB} \times m_{PB}\).
Even when the population-mean \(PB\) amplitude is zero, individual
participants have nonzero \(PB\) draws (standard deviation
\(\sigma_{PB}=2\)), and coupling the biomarker to that variance is
sufficient to contaminate the estimator. A true no-\(PB\) control
requires \(\sigma_{PB}=0\), not \(m_{PB}=0\). This refines the
Study A design: varying \(m_{PB}\) at \(c_{bm,PB}=0\) correctly
produced zero bias, but because the relevant null was
\(c_{bm,PB}=0\), not \(m_{PB}=0\).

The phase-augmented analysis provides no protection: its bias
at every \((c_{bm,PB}, m_{PB}, N)\) cell is indistinguishable
from the one-component bias. Adding phase-by-treatment
interactions to the model does not decompose the contamination
because both analyses see the biomarker-\(PB\) covariance through
the same open-label versus blinded-discontinuation contrast;
the contrast changes the weight on \(PB\) (via expectancy) but
does not isolate the \(PB\) component's contribution to the
biomarker slope. Separating that contribution requires both a
design that decouples drug exposure from belief and a
decomposition that exploits the decoupling; Study B (section 7.5)
shows that the standard hybrid design cannot achieve this
separation, whereas a balanced-placebo design can.

The \(N=70\) results are uninformative at 100 replicates (Monte
Carlo standard error \(\approx 0.06\), comparable to or larger
than the bias at moderate \(c_{bm,PB}\)). The production run
will use 1,000 replicates per cell to establish the \(N=70\)
pattern and 5,000 null replicates to assess Type I error under
contamination. Until the production run is complete, the
quantitative bias estimates from this pilot should be treated
as directionally informative but imprecise.

\end{orig}

\begin{rgt}
Lorem ipsum dolor sit amet, consectetur adipiscing elit, sed do eiusmod tempor incididunt ut labore et dolore magna aliqua.

\end{rgt}

\begin{orig}
The pilot results reframe the paper's practical guidance. The
risk of biomarker-validation failure in an aggregated N-of-1
study arises not from the presence of a strong placebo response
in the population but from the biomarker's correlation with
individual-level placebo responsiveness. In a population where
the baseline biomarker is uncorrelated with each participant's
placebo-belief response amplitude, the one-component estimator
recovers the pharmacological interaction without bias, and the
full decomposition machinery is not required. In a population
where the biomarker covaries with psychological characteristics
that predict placebo response -- optimizm, prior treatment
experience, suggestibility, expectancy at enrollment -- the
one-component analysis conflates pharmacological and
placebo-belief prediction, and the scale of that conflation
grows with the coupling strength. At \(c_{bm,PB}=0.45\), a
value that equals the pharmacological coupling used throughout
this paper, the one-component bias is approximately \(0.51\)
points on the interaction scale, representing \(23\%\) of the
true pharmacological effect. This is a practically meaningful
distortion for any precision-medicine stratification decision.

The production run of the contaminated-biomarker study, at
1,000 alternative-cell replicates and 5,000 null-cell
replicates, is required to confirm these estimates, establish
the \(N=70\) pattern, and characterize Type I error under
contamination. Results are pending.

\end{orig}

\subsection{Study B: recovery under a belief-decoupling design}\label{study-b-recovery-under-a-belief-decoupling-design}

\begin{rgt}
Lorem ipsum dolor sit amet, consectetur adipiscing elit, sed do eiusmod tempor incididunt ut labore et dolore magna aliqua.

\end{rgt}

\begin{orig}
Study B asks whether the contamination demonstrated in section 7.4 can
be removed by decomposition. The answer is conditional on the trial
design. The mechanism is that, in the multivariate normal (MVN) DGP, the
biomarker-\(BR\) coupling acts only on drug (and so loads on the drug-state
contrast \(D_{bc}\)), whereas the biomarker-\(PB\) coupling scales with the
expectancy \(\eta\) (the placebo component's amplitude is proportional to
\(\eta\), so its biomarker-coupled part loads on \(\eta\)). A decomposition
that interacts the biomarker with both \(D_{bc}\) and \(\eta\) can therefore
isolate \(\beta_{bm}^{BR}\) on \(bm{:}D_{bc}\), but only if \(D_{bc}\) and
\(\eta\) vary independently in the design. In the standard hybrid
open-label-plus-blinded-discontinuation design they do not: on-drug time
and belief accumulate together (their analysis-stage correlation is
about \(+0.6\)), and we verified that the belief-covariate decomposition,
a blinded-stratum drug contrast, and a Gompertz-basis component
decomposition all fail to recover \(\beta_{bm}^{BR}\) on that design,
performing no better than, and sometimes worse than, the one-component
analysis. The remedy is a design that breaks the collinearity. A
balanced-placebo (open-hidden) design\textsuperscript{\citeproc{ref-colloca2004}{51},\citeproc{ref-rohsenow1981}{52}}
interposes a covert on-drug phase, in which the drug is administered
without the participant's knowledge (\(\eta=0\)), and an open-placebo
phase, in which the participant is off drug but believes treatment
continues (\(\eta=1\)); the first supplies drug variation at fixed belief
and the second belief variation at fixed drug exposure, reducing the
\(D_{bc}\)-\(\eta\) correlation to about \(-0.66\).

\end{orig}

\begin{rgt}
Lorem ipsum dolor sit amet, consectetur adipiscing elit, sed do eiusmod tempor incididunt ut labore et dolore magna aliqua.

\end{rgt}

\begin{orig}
The production run (1\{,\}000 replicates per cell, balanced-placebo design
with three open-label, six covert on-drug, and six open-placebo
timepoints, \(N \in \{70, 300\}\), \(c_{bm,PB} \in \{0, 0.1, 0.2, 0.3\}\))
confirms the recovery at full resolution. At \(N=300\) the mean estimated
\(\hat\beta_{bm:D_{bc}}\) from the decomposition tracks the pharmacological
target \(-2.25\) across the entire contamination range, with bias
\(-0.022\), \(+0.008\), \(+0.021\), and \(-0.026\) at \(c_{bm,PB}=0, 0.1, 0.2,
0.3\) respectively, every value within one MCSE
(about \(0.02\)) of zero. The one-component estimator on the same data is
unbiased only at \(c_{bm,PB}=0\) and develops a monotone bias of \(+0.16\),
\(+0.34\), and \(+0.48\) as the biomarker-\(PB\) correlation grows. Calibration
tracks the bias: the decomposition holds near-nominal interval coverage
(\(0.94\) to \(0.96\)), whereas one-component coverage degrades to \(0.86\) at
\(c_{bm,PB}=0.3\). The \(N=70\) cells show the same pattern with the
decomposition unbiased (absolute bias \(\le 0.08\)) and lower power
(about \(0.66\) versus \(0.93\) at \(N=300\)), the precision cost of the
smaller sample and the extra belief-covariate terms. Convergence was
complete in every cell. Per-replicate output and the cell-level summary
are at
\texttt{analysis/data/derived\_data/component-decomposition/study-b-balanced/}.
The contamination range is bounded above by positive-definiteness: at
\(c_{bm,PB}=0.45\) the biomarker would have to correlate \(0.45\) with both
\(BR\) and \(PB\), which the covariance cannot accommodate (minimum
eigenvalue \(-0.35\)), so we restrict the production grid to the feasible
range \(c_{bm,PB} \le 0.30\).

\end{orig}

\begin{rgt}
Lorem ipsum dolor sit amet, consectetur adipiscing elit, sed do eiusmod tempor incididunt ut labore et dolore magna aliqua.

\end{rgt}

\begin{orig}
The result sharpens the paper's central claim in an important way. The
inferential failure of section 7.4 is real and the decomposition does
remove it, but the decomposition is not self-sufficient: on the standard
hybrid design it does not recover the pharmacological slope, because the
design does not supply independent variation in drug exposure and belief.
The enabling condition is a design that decouples the two, of which the
balanced-placebo and open-hidden paradigms\textsuperscript{\citeproc{ref-colloca2004}{51},\citeproc{ref-rohsenow1981}{52}}
are the established instances. The practical recommendation that follows
is specific: a biomarker-validation trial in which the candidate
biomarker may correlate with placebo responsiveness should be fielded as
a balanced-placebo or open-hidden design rather than the standard hybrid
design, because only the former makes the pharmacological
biomarker-treatment interaction separately identifiable. Three
questions remain open. Recovery beyond the positive-definite contamination
range is not characterized, since the MVN architecture cannot represent
it; the participant-level component recovery that motivated the original
identifiability aims of Study B is not addressed by the slope-level
analysis here; and whether the same recovery holds under the mean-moderation architecture
(mean-moderation) coupling, where the biomarker shifts component means
rather than correlating with component draws, remains to be confirmed.

\end{orig}

\section{Application to predictive-biomarker validation}\label{application-to-predictive-biomarker-validation}

\begin{rgt}
Lorem ipsum dolor sit amet, consectetur adipiscing elit, sed do eiusmod tempor incididunt ut labore et dolore magna aliqua.

\end{rgt}

\begin{orig}
The motivating clinical question for the trial designs in this
project is whether a baseline biomarker predicts treatment
response, and specifically whether the biomarker predicts the
pharmacological component of treatment response. This is the
precision-medicine question, formalized in the PATH statement\textsuperscript{\citeproc{ref-kent2020path}{14}}. Whether it depends on the BR-PB-TV decomposition
is itself conditional: by the identity of section 6.1, the one-component
biomarker slope equals the pharmacological slope when the biomarker
is uncorrelated with \(PB\) and \(TV\), and departs from it only by the
biomarker's covariance with those components. The decomposition is
therefore the mechanism the question requires precisely in the
regime where that covariance is non-negligible.

\end{orig}

\begin{rgt}
Lorem ipsum dolor sit amet, consectetur adipiscing elit, sed do eiusmod tempor incididunt ut labore et dolore magna aliqua.

\end{rgt}

\begin{orig}
A direct regression of total response on baseline biomarker

\begin{verbatim}
total_response = beta_0 + beta_bm * biomarker + error
\end{verbatim}

estimates the biomarker's correlation with whatever the one-component
total contains. If \(PB\) varies systematically across the
biomarker distribution, because the biomarker is correlated
with optimizm, suggestibility, or treatment-seeking history,
then \(\hat{\beta}_{bm}\) contains a placebo-prediction component
inseparable from the pharmacological-prediction component. The
resulting biomarker, applied prospectively, would mis-stratify
patients: those flagged as `high responders' would include both
true high-\(BR\) responders and high-\(PB\) responders, and the
latter would not benefit from prescription decisions made on a
pharmacological basis.

\end{orig}

\begin{rgt}
Lorem ipsum dolor sit amet, consectetur adipiscing elit, sed do eiusmod tempor incididunt ut labore et dolore magna aliqua.

\end{rgt}

\begin{orig}
The right approach is to regress each component on the
biomarker:

\begin{align*}
BR &= \beta_0^{BR} + \beta_{bm}^{BR} \cdot \mathrm{biomarker} +
\varepsilon^{BR},\\
PB &= \beta_0^{PB} + \beta_{bm}^{PB} \cdot \mathrm{biomarker} +
\varepsilon^{PB},\\
TV &= \beta_0^{TV} + \beta_{bm}^{TV} \cdot \mathrm{biomarker} +
\varepsilon^{TV}.
\end{align*}

A biomarker is a useful pharmacological predictor if and only
if \(\hat{\beta}_{bm}^{BR}\) is meaningfully non-zero. Significant
\(\hat{\beta}_{bm}^{PB}\) or \(\hat{\beta}_{bm}^{TV}\) are
diagnostic findings about the biomarker, indicating that it is
partly psychological or correlates with natural-history
characteristics, but they are not the validation of
pharmacological prediction. A biomarker that correlates only
with \(PB\) would be valuable for selecting placebo responders in
run-in designs but worthless for predicting pharmacological
efficacy; a biomarker that correlates only with \(TV\) would be
a prognostic marker, not a predictive one.

\end{orig}

\begin{rgt}
Lorem ipsum dolor sit amet, consectetur adipiscing elit, sed do eiusmod tempor incididunt ut labore et dolore magna aliqua.

\end{rgt}

\begin{orig}
The decomposition is therefore the mechanism by which the
precision-medicine question is posed whenever the candidate
biomarker may correlate with \(PB\) or \(TV\). Where the biomarker can
be assumed orthogonal to the non-pharmacological components, the
identity of section 6.1 shows the one-component analysis already targets
the biomarker-by-\(BR\) estimand, and the simulation of section 7.3
confirms it recovers that estimand without bias. The phase-augmented
linear-mixed analysis introduced in section 5 retains a useful
subset of the decomposition machinery at small \(N\) (the
\texttt{bm:Dbc:phase} interaction tests whether the
biomarker-treatment effect attenuates under blinding, the
practical proxy for biomarker-by-\(PB\) contamination), but Study A
shows this proxy earns its degrees of freedom only when such
contamination is present: under an uncontaminated biomarker it
reduced no bias and lost power. The full decomposition recovers the
component-specific regressions explicitly when \(N\) permits and when
the scientific question warrants the cost.

\end{orig}

\section{Discussion}\label{discussion}

\begin{rgt}
Lorem ipsum dolor sit amet, consectetur adipiscing elit, sed do eiusmod tempor incididunt ut labore et dolore magna aliqua.

\end{rgt}

\begin{orig}
The three-component decomposition is the formal expression of
the claim that `treatment response' in an aggregated N-of-1
trial is not a single quantity but a sum of three causally
distinct contributions, each with its own pharmacological,
psychological, or epidemiological interpretation. What the present
work establishes, analytically in section 6.1 and empirically in
section 7.3, is that the inferential value of acting on this
decomposition is conditional rather than universal. The identity
shows that the one-component estimator is displaced from its
pharmacological target only by the components' loading onto the
contrast used, and, for a biomarker slope, only by the biomarker's
covariance with \(PB\) and \(TV\). Where those terms are zero, the
design-contrast one-component analysis is already unbiased, and the
production simulation confirms it: with a biomarker coupled to \(BR\)
alone, the one-component biomarker-treatment estimate is
unbiased across the whole \(PB\)-\(TV\) grid and the phase-augmented
analysis is dominated. The framework's contribution is therefore
best read as a map of when decomposition changes inference, not as
a blanket prescription to decompose. It makes the trial-design
conditions explicit, identifies the attribution and
contaminated-biomarker estimands for which decomposition is needed,
and clarifies the target estimand for biomarker validation (the
biomarker-by-\(BR\) association, not the one-component total).

\end{orig}

\begin{rgt}
Lorem ipsum dolor sit amet, consectetur adipiscing elit, sed do eiusmod tempor incididunt ut labore et dolore magna aliqua.

\end{rgt}

\begin{orig}
Four open questions remain and warrant acknowledgement. First,
the simulation study described in section 7 is forthcoming; the
present manuscript reports the framework rather than its
empirical evaluation, and the quantitative bias-versus-power
trade-offs across analysis strategies are conjectures until
that work is run. Second, the \(\eta(\text{phase})\) multipliers
used to parameterize the expectancy modulation of \(PB\) are
themselves modeling assumptions, and a sensitivity analysis
varying \(\eta\) across plausible values is mandatory before any
clinical conclusion is drawn from the framework's outputs.
Third, the framework as specified treats the three components
as monotone Gompertz trajectories; cyclical or biphasic \(TV\)
patterns require either flexible \(TV\) specifications or explicit
covariates beyond the Gompertz form, and the boundary between
`flexible-Gompertz' and `covariate-augmented' specifications is
not yet sharply drawn. Fourth, the participant-level component
estimates the framework targets inherit the identifiability
caveat that\textsuperscript{\citeproc{ref-senn2018}{54}} and\textsuperscript{\citeproc{ref-senn2004}{49}} raise for individual
treatment effects more generally: within-participant component
recovery rests on separating a true participant-by-component
signal from participant-by-period and participant-by-phase
variation, and this separation is credible only when the design
supplies repeated phase contrasts and the cross-component
covariance structure is not pathological. The framework should
be read as estimating these quantities under those assumptions,
not as dissolving the well-known difficulty of estimating
individual effects from finite within-person data.

\end{orig}

\begin{rgt}
Lorem ipsum dolor sit amet, consectetur adipiscing elit, sed do eiusmod tempor incididunt ut labore et dolore magna aliqua.

\end{rgt}

\begin{orig}
The companion pedagogical document at
\texttt{docs/24-component-decomposition-pedagogy.md} develops
the framework with worked examples and an extended intuition
narrative for readers approaching the material for the first
time. The companion methodological evaluation of the Gompertz
family at \texttt{analysis/report/07-gompertz-evaluation/}
quantifies the sensitivity of the framework's inference to the
specific parametric form chosen for the component trajectories.

\end{orig}

\begin{rgt}
Lorem ipsum dolor sit amet, consectetur adipiscing elit, sed do eiusmod tempor incididunt ut labore et dolore magna aliqua.

\end{rgt}

\begin{orig}
A final point concerns the scope of the decomposition, where it is
essential to distinguish the existence of the components from their
identifiability. The claim that an observed longitudinal treatment
response is the sum of a pharmacological, a belief-driven, and a
natural-history component is a statement about the data-generating
process of any repeated-measures treatment study, and in that sense the
decomposition is not specific to N-of-1 trials. Identifiability, by
contrast, is design-specific, and the distinction between
population-level and participant-level recovery is decisive for the
biomarker application. A standard parallel-group placebo-controlled
trial supplies a between-arm placebo contrast that identifies the
population-mean belief response but not a participant-level \(PB\), and it
contains no within-subject contrast in which belief varies independently
of pharmacology; the participant-level biomarker-by-\(BR\) slope that the
precision-medicine question targets is therefore not identified from
such a design, no matter how large it is. The designs that do supply the
required within-subject belief contrast are the blinded or open-hidden
discontinuation and randomized-withdrawal designs\textsuperscript{\citeproc{ref-colloca2004}{51},\citeproc{ref-rohsenow1981}{52}}, the crossover with on/off periods, and
the hybrid N-of-1 design developed here; the pharmacometric
disease-progression-plus-placebo models from which the framework
descends\textsuperscript{\citeproc{ref-holford2015}{1}--\citeproc{ref-chen2021}{3}} operate at the group
level and recover population-mean components only. We therefore advocate
the decomposition as a default lens for longitudinal designs that supply
a within-subject belief-by-exposure contrast, with the aggregated N-of-1
case the most demanding instance because it must recover the components
from a single participant's trajectory. We do not claim participant-level
identifiability from parallel-group placebo control alone, and where only
a between-arm placebo contrast is available the framework's claims are
restricted to the population mean.

\end{orig}

\section{Conclusions}\label{conclusions}

\begin{rgt}
Lorem ipsum dolor sit amet, consectetur adipiscing elit, sed do eiusmod tempor incididunt ut labore et dolore magna aliqua.

\end{rgt}

\begin{orig}
In conclusion, three points are to be emphasized. First,
treatment response in aggregated N-of-1 clinical trials is the
sum of three causally distinct components: pharmacological
action (\(BR\)), placebo-belief response (\(PB\)), and time-variant
natural history (\(TV\)). Each component has its own biological
mechanism, its own identifiability conditions, and its own
clinical interpretation. Whether conflating them biases a given
analysis is, however, governed by the omitted-variable-bias
identity of section 6.1, not by the raw magnitude of \(PB\) and
\(TV\): the one-component pre-post treatment effect is displaced by the
\(PB\) and \(TV\) contributions, whereas the biomarker-treatment
slope is displaced only by the biomarker's covariance with those
components. Study A bears this out: varying \(PB\) and \(TV\) magnitude
with a biomarker coupled to \(BR\) alone left the one-component
estimate unbiased across the entire grid.

\end{orig}

\begin{rgt}
Lorem ipsum dolor sit amet, consectetur adipiscing elit, sed do eiusmod tempor incididunt ut labore et dolore magna aliqua.

\end{rgt}

\begin{orig}
Second, the operative requirement is the design contrast, not the
parametric analysis model. The trial design must include the
drug-on/off and belief contrasts that identify the components, and
the one-component analysis already exploits the drug-on/off
contrast through its \(D_{bc}\) term. Adding parametric component
structure on top is not free: under an uncontaminated biomarker the
phase-augmented analysis reduced no bias and lost substantial
power, so it should not be adopted as a default. Its value, and
that of the full component-matched fit, is confined to the
estimands where the extra structure does work, namely the
attribution of response to pharmacology versus belief and the
analysis of biomarkers that may themselves correlate with \(PB\) or
\(TV\).

\end{orig}

\begin{rgt}
Lorem ipsum dolor sit amet, consectetur adipiscing elit, sed do eiusmod tempor incididunt ut labore et dolore magna aliqua.

\end{rgt}

\begin{orig}
Third, for predictive-biomarker validation the decomposition earns
its place precisely when the candidate biomarker may correlate with
\(PB\) or \(TV\). A biomarker that predicts the one-component total is a
pharmacological predictor only when it is uncorrelated with the
non-pharmacological components; where that cannot be assumed, the
one-component slope mixes pharmacological and non-pharmacological
prediction, and the regulatory and prescribing decisions that
motivate the trial require the separation the decomposition
supplies. The contaminated-biomarker study (section 7.4) demonstrates
this failure directly, and Study B (section 7.5) establishes that the
bias is recoverable: a belief-covariate decomposition returns the
unbiased pharmacological slope across the feasible contamination range,
but only on a balanced-placebo design that decouples drug exposure from
belief, not on the standard hybrid design. The decisive practical
lesson is therefore as much about design as about analysis.

\end{orig}

\section{Conflict of interest}\label{conflict-of-interest}

The authors declare no conflicts of interest.

\section{Funding}\label{funding}

This work received no specific funding.

\section{Data availability statement}\label{data-availability-statement}

The data and code that support the findings of this study are
openly available in the pmsimstats-ng project repository.
Per-replicate Monte Carlo outputs and ADEMP-compliant cell-level
summaries are versioned under \texttt{analysis/data/};
simulation drivers are at \texttt{analysis/scripts/}. Exact
paths for this manuscript are listed in the Reproducibility
section below.

\section{Reproducibility}\label{reproducibility}

\begin{rgt}
Lorem ipsum dolor sit amet, consectetur adipiscing elit, sed do eiusmod tempor incididunt ut labore et dolore magna aliqua.

\end{rgt}

\begin{orig}
This manuscript was rendered with R 4.4.x inside the project's
zzcollab Docker container (image hash recorded in
\texttt{Dockerfile}; package set captured in \texttt{renv.lock}).
Principal packages used by the simulation pipeline are
\texttt{nlme} (continuous-time linear-mixed analysis with
\texttt{corCAR1} residual correlation), \texttt{parallel}
(8-core \texttt{mclapply} for parameter-grid sweeps),
\texttt{data.table} (the \texttt{original-extended} simulation
collection), \texttt{furrr} and \texttt{purrr} (the
\texttt{tidyverse} simulation collection), and
\texttt{rmarkdown} plus \texttt{bookdown} for the manuscript
build.

\end{orig}

\begin{rgt}
Lorem ipsum dolor sit amet, consectetur adipiscing elit, sed do eiusmod tempor incididunt ut labore et dolore magna aliqua.

\end{rgt}

\begin{orig}
Each simulation driver sets \texttt{set.seed(20260507)} at the
top of the replicate loop and uses
\texttt{parallel::mc.set.seed} inside \texttt{mclapply}.
Per-replicate seeds derive from the master seed by
\texttt{+ rep\_idx}.

\end{orig}

\begin{rgt}
Lorem ipsum dolor sit amet, consectetur adipiscing elit, sed do eiusmod tempor incididunt ut labore et dolore magna aliqua.

\end{rgt}

\begin{orig}
Per-replicate simulation outputs are checkpointed at
\texttt{analysis/data/quick-sim/06-decomp-replicates.rds};
Morris ADEMP cell-level summaries at the companion
\texttt{06-decomp-summary.txt}. Driver scripts are under
\texttt{analysis/scripts/component-decomposition/}. The
pre-registered ADEMP plan is at
\texttt{analysis/scripts/component-decomposition/00-ademp-pre-registration.md}.
The manuscript source, paper-specific bibliography
(\texttt{references.bib}), and SIM CSL file are at
\texttt{analysis/report/06-component-decomposition/}.

\end{orig}

\section{References}\label{references}

\protect\phantomsection\label{refs}
\begin{CSLReferences}{0}{1}
\bibitem[\citeproctext]{ref-holford2015}
\CSLLeftMargin{1. }%
\CSLRightInline{Holford N. Clinical pharmacology = disease progression + drug action. \emph{British Journal of Clinical Pharmacology}. 2015;79(1):18-27. doi:\href{https://doi.org/10.1111/bcp.12170}{10.1111/bcp.12170}}

\bibitem[\citeproctext]{ref-gomeni2009}
\CSLLeftMargin{2. }%
\CSLRightInline{Gomeni R, Lavergne A, Merlo-Pich E. Modelling placebo response in depression trials using a longitudinal model with informative dropout. \emph{European Journal of Pharmaceutical Sciences}. 2009;36(1):4-10. doi:\href{https://doi.org/10.1016/j.ejps.2008.10.025}{10.1016/j.ejps.2008.10.025}}

\bibitem[\citeproctext]{ref-chen2021}
\CSLLeftMargin{3. }%
\CSLRightInline{Chen C, Jönsson S, Yang S, Plan EL, Karlsson MO. Detecting placebo and drug effects on {Parkinson's} disease symptoms by longitudinal item-score models. \emph{CPT: Pharmacometrics \& Systems Pharmacology}. 2021;10(4):309-317. doi:\href{https://doi.org/10.1002/psp4.12601}{10.1002/psp4.12601}}

\bibitem[\citeproctext]{ref-muthen2009}
\CSLLeftMargin{4. }%
\CSLRightInline{Muthén B, Brown HC. Estimating drug effects in the presence of placebo response: Causal inference using growth mixture modeling. \emph{Statistics in Medicine}. 2009;28(27):3363-3385. doi:\href{https://doi.org/10.1002/sim.3721}{10.1002/sim.3721}}

\bibitem[\citeproctext]{ref-kaptchuk2008components}
\CSLLeftMargin{5. }%
\CSLRightInline{{Kaptchuk TJ, Kelley JM, Conboy LA, et al.} Components of placebo effect: Randomised controlled trial in patients with irritable bowel syndrome. \emph{BMJ}. 2008;336(7651):999-1003. doi:\href{https://doi.org/10.1136/bmj.39524.439618.25}{10.1136/bmj.39524.439618.25}}

\bibitem[\citeproctext]{ref-kirsch2008}
\CSLLeftMargin{6. }%
\CSLRightInline{Kirsch I, Deacon BJ, Huedo-Medina TB, Scoboria A, Moore TJ, Johnson BT. Initial severity and antidepressant benefits: A meta-analysis of data submitted to the {Food and Drug Administration}. \emph{PLoS Medicine}. 2008;5(2):e45. doi:\href{https://doi.org/10.1371/journal.pmed.0050045}{10.1371/journal.pmed.0050045}}

\bibitem[\citeproctext]{ref-locher2015}
\CSLLeftMargin{7. }%
\CSLRightInline{{Locher C, Kossowsky J, Gaab J, Kirsch I, et al.} Moderation of antidepressant and placebo outcomes by baseline severity in late-life depression: A systematic review and meta-analysis. \emph{Journal of Affective Disorders}. 2015;181:50-60. doi:\href{https://doi.org/10.1016/j.jad.2015.03.062}{10.1016/j.jad.2015.03.062}}

\bibitem[\citeproctext]{ref-sugarman2014}
\CSLLeftMargin{8. }%
\CSLRightInline{Sugarman MA, Loree AM, Baltes BB, Grekin ER, Kirsch I. The efficacy of paroxetine and placebo in treating anxiety and depression. \emph{PLoS ONE}. 2014;9(8):e106337. doi:\href{https://doi.org/10.1371/journal.pone.0106337}{10.1371/journal.pone.0106337}}

\bibitem[\citeproctext]{ref-hall2012comt}
\CSLLeftMargin{9. }%
\CSLRightInline{Hall KT, Lembo AJ, Kirsch I, et al. Catechol-{O}-methyltransferase val158met polymorphism predicts placebo effect in irritable bowel syndrome. \emph{PLoS ONE}. 2012;7(10):e48135. doi:\href{https://doi.org/10.1371/journal.pone.0048135}{10.1371/journal.pone.0048135}}

\bibitem[\citeproctext]{ref-wang2022comt}
\CSLLeftMargin{10. }%
\CSLRightInline{{Wang RS, Lembo AJ, Kaptchuk TJ, Hall KT, et al.} Genomic effects associated with response to placebo treatment in a randomized trial of irritable bowel syndrome. \emph{Frontiers in Pain Research}. 2022;2:775386. doi:\href{https://doi.org/10.3389/fpain.2021.775386}{10.3389/fpain.2021.775386}}

\bibitem[\citeproctext]{ref-hrobjartsson2001}
\CSLLeftMargin{11. }%
\CSLRightInline{Hróbjartsson A, Gøtzsche PC. Is the placebo powerless? An analysis of clinical trials comparing placebo with no treatment. \emph{New England Journal of Medicine}. 2001;344(21):1594-1602. doi:\href{https://doi.org/10.1056/NEJM200105243442106}{10.1056/NEJM200105243442106}}

\bibitem[\citeproctext]{ref-hrobjartsson2004}
\CSLLeftMargin{12. }%
\CSLRightInline{Hróbjartsson A, Gøtzsche PC. Is the placebo powerless? Update of a systematic review with 52 new randomized trials. \emph{Journal of Internal Medicine}. 2004;256(2):91-100. doi:\href{https://doi.org/10.1111/j.1365-2796.2004.01355.x}{10.1111/j.1365-2796.2004.01355.x}}

\bibitem[\citeproctext]{ref-hrobjartsson2010}
\CSLLeftMargin{13. }%
\CSLRightInline{Hróbjartsson A, Gøtzsche PC. Placebo interventions for all clinical conditions. \emph{Cochrane Database of Systematic Reviews}. 2010;(1):CD003974. doi:\href{https://doi.org/10.1002/14651858.CD003974.pub3}{10.1002/14651858.CD003974.pub3}}

\bibitem[\citeproctext]{ref-kent2020path}
\CSLLeftMargin{14. }%
\CSLRightInline{{Kent DM, Paulus JK, Klaveren D van, et al.} The {P}redictive {A}pproaches to {T}reatment effect {H}eterogeneity ({PATH}) {S}tatement. \emph{Annals of Internal Medicine}. 2020;172(1):35-45. doi:\href{https://doi.org/10.7326/M18-3667}{10.7326/M18-3667}}

\bibitem[\citeproctext]{ref-kent2020pathee}
\CSLLeftMargin{15. }%
\CSLRightInline{{Kent DM, Klaveren D van, Paulus JK, et al.} The {PATH} statement: Explanation and elaboration. \emph{Annals of Internal Medicine}. 2020;172(1):W1-W25. doi:\href{https://doi.org/10.7326/M18-3668}{10.7326/M18-3668}}

\bibitem[\citeproctext]{ref-kent2018bmj}
\CSLLeftMargin{16. }%
\CSLRightInline{Kent DM, Steyerberg E, Klaveren D van. Personalized evidence based medicine: Predictive approaches to heterogeneous treatment effects. \emph{BMJ}. 2018;363:k4245. doi:\href{https://doi.org/10.1136/bmj.k4245}{10.1136/bmj.k4245}}

\bibitem[\citeproctext]{ref-rekkas2020scoping}
\CSLLeftMargin{17. }%
\CSLRightInline{{Rekkas A, Paulus JK, Raman G, et al.} Predictive approaches to heterogeneous treatment effects: A scoping review. \emph{BMC Medical Research Methodology}. 2020;20(1):264. doi:\href{https://doi.org/10.1186/s12874-020-01145-1}{10.1186/s12874-020-01145-1}}

\bibitem[\citeproctext]{ref-guyatt1986}
\CSLLeftMargin{18. }%
\CSLRightInline{Guyatt G, Sackett D, Taylor DW, Chong J, Roberts R, Pugsley S. Determining optimal therapy: Randomized trials in individual patients. \emph{New England Journal of Medicine}. 1986;314(14):889-892.}

\bibitem[\citeproctext]{ref-zucker1997}
\CSLLeftMargin{19. }%
\CSLRightInline{Zucker DR, Schmid CH, McIntosh MW, D'Agostino RB, Selker HP, Lau J. Combining single patient ({N}-of-1) trials to estimate population treatment effects and to evaluate individual patient responses to treatment. \emph{Journal of Clinical Epidemiology}. 1997;50(4):401-410. doi:\href{https://doi.org/10.1016/S0895-4356(96)00429-5}{10.1016/S0895-4356(96)00429-5}}

\bibitem[\citeproctext]{ref-zucker2010}
\CSLLeftMargin{20. }%
\CSLRightInline{Zucker DR, Ruthazer R, Schmid CH. Individual ({N}-of-1) trials can be combined to give population comparative treatment effect estimates: Methodologic considerations. \emph{Journal of Clinical Epidemiology}. 2010;63(12):1312-1323. doi:\href{https://doi.org/10.1016/j.jclinepi.2010.04.020}{10.1016/j.jclinepi.2010.04.020}}

\bibitem[\citeproctext]{ref-lillie2011}
\CSLLeftMargin{21. }%
\CSLRightInline{Lillie EO, Patay B, Diamant J, Issell B, Topol EJ, Schork NJ. The n-of-1 clinical trial: The ultimate strategy for individualizing medicine? \emph{Personalized Medicine}. 2011;8(2):161-173. doi:\href{https://doi.org/10.2217/pme.11.7}{10.2217/pme.11.7}}

\bibitem[\citeproctext]{ref-duan2013}
\CSLLeftMargin{22. }%
\CSLRightInline{Duan N, Kravitz RL, Schmid CH. Single-patient ({n}-of-1) trials: A pragmatic clinical decision methodology for patient-centered comparative effectiveness research. \emph{Journal of Clinical Epidemiology}. 2013;66(8 Suppl):S21-S28. doi:\href{https://doi.org/10.1016/j.jclinepi.2013.04.006}{10.1016/j.jclinepi.2013.04.006}}

\bibitem[\citeproctext]{ref-gabler2011}
\CSLLeftMargin{23. }%
\CSLRightInline{Gabler NB, Duan N, Vohra S, Kravitz RL. {N}-of-1 trials in the medical literature: A systematic review. \emph{Medical Care}. 2011;49(8):761-768. doi:\href{https://doi.org/10.1097/MLR.0b013e318215d90d}{10.1097/MLR.0b013e318215d90d}}

\bibitem[\citeproctext]{ref-schork2017nutrition}
\CSLLeftMargin{24. }%
\CSLRightInline{Schork NJ, Goetz LH. Single-subject studies in translational nutrition research. \emph{Annual Review of Nutrition}. 2017;37:395-422. doi:\href{https://doi.org/10.1146/annurev-nutr-071816-064717}{10.1146/annurev-nutr-071816-064717}}

\bibitem[\citeproctext]{ref-schorkRandomizedClinicalTrials2018}
\CSLLeftMargin{25. }%
\CSLRightInline{Schork NJ. Randomized clinical trials and personalized medicine: A commentary on {Deaton} and {Cartwright}. \emph{Social Science and Medicine}. 2018;210:71-73. doi:\href{https://doi.org/10.1016/j.socscimed.2018.04.033}{10.1016/j.socscimed.2018.04.033}}

\bibitem[\citeproctext]{ref-schork2023precision}
\CSLLeftMargin{26. }%
\CSLRightInline{Schork NJ, Beaulieu-Jones B, Liang WS, Smalley S, Goetz LH. Exploring human biology with {N}-of-1 clinical trials. \emph{Cambridge Prisms: Precision Medicine}. 2023;1:e12. doi:\href{https://doi.org/10.1017/pcm.2022.15}{10.1017/pcm.2022.15}}

\bibitem[\citeproctext]{ref-porcino2022}
\CSLLeftMargin{27. }%
\CSLRightInline{Porcino A, Vohra S. {N}-of-1 trials, their reporting guidelines, and the advancement of open science principles. \emph{Harvard Data Science Review}. 2022;4(SI3). doi:\href{https://doi.org/10.1162/99608f92.a65a257a}{10.1162/99608f92.a65a257a}}

\bibitem[\citeproctext]{ref-vohraCONSORTExtensionReporting2015}
\CSLLeftMargin{28. }%
\CSLRightInline{{Vohra S, Shamseer L, Sampson M, et al.} {CONSORT} extension for reporting {N}-of-1 trials ({CENT}) 2015 statement. \emph{Journal of Clinical Epidemiology}. 2015;76:9-17. doi:\href{https://doi.org/10.1016/j.jclinepi.2015.05.004}{10.1016/j.jclinepi.2015.05.004}}

\bibitem[\citeproctext]{ref-shamseer2015cent}
\CSLLeftMargin{29. }%
\CSLRightInline{{Shamseer L, Sampson M, Bukutu C, Schmid CH, et al.} {CONSORT} extension for reporting {N}-of-1 trials ({CENT}) 2015: Explanation and elaboration. \emph{Journal of Clinical Epidemiology}. 2015;76:18-46. doi:\href{https://doi.org/10.1016/j.jclinepi.2015.05.018}{10.1016/j.jclinepi.2015.05.018}}

\bibitem[\citeproctext]{ref-li2016cent}
\CSLLeftMargin{30. }%
\CSLRightInline{Li J, Gao W, Punja S, et al. Reporting quality of {N}-of-1 trials published between 1985 and 2013: A systematic review. \emph{Journal of Clinical Epidemiology}. 2016;76:57-64. doi:\href{https://doi.org/10.1016/j.jclinepi.2015.11.016}{10.1016/j.jclinepi.2015.11.016}}

\bibitem[\citeproctext]{ref-liao2023}
\CSLLeftMargin{31. }%
\CSLRightInline{Liao Z, Qian M, Kronish IM, Cheung YK. Analysis of {N}-of-1 trials using {B}ayesian distributed lag model with autocorrelated errors. \emph{Statistics in Medicine}. 2023;42(13):2044-2060. doi:\href{https://doi.org/10.1002/sim.9676}{10.1002/sim.9676}}

\bibitem[\citeproctext]{ref-blackstonComparisonAggregatedNof12019}
\CSLLeftMargin{32. }%
\CSLRightInline{Blackston JW, Chapple AG, McGree JM, McDonald S, Nikles J. Comparison of aggregated {N}-of-1 trials with parallel and crossover randomized controlled trials using simulation studies. \emph{Healthcare}. 2019;7(4):137. doi:\href{https://doi.org/10.3390/healthcare7040137}{10.3390/healthcare7040137}}

\bibitem[\citeproctext]{ref-carrozzo2025}
\CSLLeftMargin{33. }%
\CSLRightInline{{Carrozzo AE, Zimmermann G, Bathke AC, et al.} Two-arm crossover randomized controlled trial versus meta-analysis of {N}-of-1 studies: Comparison of statistical efficiency. \emph{Biometrical Journal}. 2025;67(2):e70045. doi:\href{https://doi.org/10.1002/bimj.70045}{10.1002/bimj.70045}}

\bibitem[\citeproctext]{ref-hendrickson2020}
\CSLLeftMargin{34. }%
\CSLRightInline{Hendrickson RC, Thomas RG, Schork NJ, Raskind MA. Optimizing aggregated {N}-of-1 trial designs for predictive biomarker validation: Statistical methods and practical considerations. \emph{Frontiers in Digital Health}. 2020;2:13. doi:\href{https://doi.org/10.3389/fdgth.2020.00013}{10.3389/fdgth.2020.00013}}

\bibitem[\citeproctext]{ref-kaptchuk2010olp}
\CSLLeftMargin{35. }%
\CSLRightInline{{Kaptchuk TJ, Friedlander E, Kelley JM, et al.} Placebos without deception: A randomized controlled trial in irritable bowel syndrome. \emph{PLoS ONE}. 2010;5(12):e15591. doi:\href{https://doi.org/10.1371/journal.pone.0015591}{10.1371/journal.pone.0015591}}

\bibitem[\citeproctext]{ref-lembo2021olp}
\CSLLeftMargin{36. }%
\CSLRightInline{{Lembo A, Kelley JM, Nee J, et al.} Open-label placebo vs double-blind placebo for irritable bowel syndrome: A randomized clinical trial. \emph{Pain}. 2021;162(9):2428-2435. doi:\href{https://doi.org/10.1097/j.pain.0000000000002234}{10.1097/j.pain.0000000000002234}}

\bibitem[\citeproctext]{ref-nurko2022}
\CSLLeftMargin{37. }%
\CSLRightInline{{Nurko S, Saps M, Kossowsky J, others, Kelley JM}. Effect of open-label placebo on children and adolescents with functional abdominal pain or irritable bowel syndrome. \emph{JAMA Pediatrics}. 2022;176(4):349-356. doi:\href{https://doi.org/10.1001/jamapediatrics.2021.5750}{10.1001/jamapediatrics.2021.5750}}

\bibitem[\citeproctext]{ref-benedetti2014}
\CSLLeftMargin{38. }%
\CSLRightInline{Benedetti F. Placebo effects: From the neurobiological paradigm to translational implications. \emph{Neuron}. 2014;84(3):623-637.}

\bibitem[\citeproctext]{ref-frisaldi2020}
\CSLLeftMargin{39. }%
\CSLRightInline{Frisaldi E, Shaibani A, Benedetti F. Understanding the mechanisms of placebo and nocebo effects. \emph{Swiss Medical Weekly}. 2020;150:w20340. doi:\href{https://doi.org/10.4414/smw.2020.20340}{10.4414/smw.2020.20340}}

\bibitem[\citeproctext]{ref-carlino2014}
\CSLLeftMargin{40. }%
\CSLRightInline{Carlino E, Frisaldi E, Benedetti F. Pain and the context. \emph{Nature Reviews Rheumatology}. 2014;10(6):348-355. doi:\href{https://doi.org/10.1038/nrrheum.2014.17}{10.1038/nrrheum.2014.17}}

\bibitem[\citeproctext]{ref-carlino2016benedetti}
\CSLLeftMargin{41. }%
\CSLRightInline{Carlino E, Piedimonte A, Benedetti F. Nature of the placebo and nocebo effect in relation to functional neurologic disorders. In: \emph{Handbook of Clinical Neurology}. Vol 139.; 2016:597-606. doi:\href{https://doi.org/10.1016/B978-0-12-801772-2.00048-5}{10.1016/B978-0-12-801772-2.00048-5}}

\bibitem[\citeproctext]{ref-wager2013nps}
\CSLLeftMargin{42. }%
\CSLRightInline{Wager TD, Atlas LY, Lindquist MA, Roy M, Woo CW, Kross E. An {fMRI}-based neurologic signature of physical pain. \emph{New England Journal of Medicine}. 2013;368(15):1388-1397. doi:\href{https://doi.org/10.1056/NEJMoa1204471}{10.1056/NEJMoa1204471}}

\bibitem[\citeproctext]{ref-zunhammer2018nps}
\CSLLeftMargin{43. }%
\CSLRightInline{Zunhammer M, Bingel U, Wager TD. Placebo effects on the neurologic pain signature: A meta-analysis of individual participant functional magnetic resonance imaging data. \emph{JAMA Neurology}. 2018;75(11):1321-1330. doi:\href{https://doi.org/10.1001/jamaneurol.2018.2017}{10.1001/jamaneurol.2018.2017}}

\bibitem[\citeproctext]{ref-raskind2013}
\CSLLeftMargin{44. }%
\CSLRightInline{{Raskind MA, Peterson K, Williams T, et al.} A trial of prazosin for combat trauma {PTSD} with nightmares in active-duty soldiers returned from {I}raq and {A}fghanistan. \emph{American Journal of Psychiatry}. 2013;170(9):1003-1010. doi:\href{https://doi.org/10.1176/appi.ajp.2013.12081133}{10.1176/appi.ajp.2013.12081133}}

\bibitem[\citeproctext]{ref-holford2011backstory}
\CSLLeftMargin{45. }%
\CSLRightInline{Holford NHG. Holford {NHG} and sheiner {LB} 'understanding the dose-effect relationship - clinical application of pharmacokinetic-pharmacodynamic models', {Clin Pharmacokin} 6:429--453 (1981) - the backstory. \emph{The AAPS Journal}. 2011;13(4):662-664. doi:\href{https://doi.org/10.1208/s12248-011-9306-5}{10.1208/s12248-011-9306-5}}

\bibitem[\citeproctext]{ref-tilburt2008}
\CSLLeftMargin{46. }%
\CSLRightInline{Tilburt JC, Emanuel EJ, Kaptchuk TJ, Curlin FA, Miller FG. Prescribing 'placebo treatments': Results of national survey of {US} internists and rheumatologists. \emph{BMJ}. 2008;337:a1938. doi:\href{https://doi.org/10.1136/bmj.a1938}{10.1136/bmj.a1938}}

\bibitem[\citeproctext]{ref-finniss2010}
\CSLLeftMargin{47. }%
\CSLRightInline{Finniss DG, Kaptchuk TJ, Miller F, Benedetti F. Biological, clinical, and ethical advances of placebo effects. \emph{The Lancet}. 2010;375(9715):686-695. doi:\href{https://doi.org/10.1016/S0140-6736(09)61706-2}{10.1016/S0140-6736(09)61706-2}}

\bibitem[\citeproctext]{ref-kirsch1985}
\CSLLeftMargin{48. }%
\CSLRightInline{Kirsch I. Response expectancy as a determinant of experience and behavior. \emph{American Psychologist}. 1985;40(11):1189-1202. doi:\href{https://doi.org/10.1037/0003-066X.40.11.1189}{10.1037/0003-066X.40.11.1189}}

\bibitem[\citeproctext]{ref-senn2004}
\CSLLeftMargin{49. }%
\CSLRightInline{Senn S. Controversies concerning randomization and additivity in clinical trials. \emph{Statistics in Medicine}. 2004;23(24):3729-3753. doi:\href{https://doi.org/10.1002/sim.2074}{10.1002/sim.2074}}

\bibitem[\citeproctext]{ref-bingel2011}
\CSLLeftMargin{50. }%
\CSLRightInline{Bingel U, Wanigasekera V, Wiech K, et al. The effect of treatment expectation on drug efficacy: Imaging the analgesic benefit of the opioid remifentanil. \emph{Science Translational Medicine}. 2011;3(70):70ra14. doi:\href{https://doi.org/10.1126/scitranslmed.3001244}{10.1126/scitranslmed.3001244}}

\bibitem[\citeproctext]{ref-colloca2004}
\CSLLeftMargin{51. }%
\CSLRightInline{Colloca L, Lopiano L, Lanotte M, Benedetti F. Overt versus covert treatment for pain, anxiety, and {Parkinson's} disease. \emph{The Lancet Neurology}. 2004;3(11):679-684.}

\bibitem[\citeproctext]{ref-rohsenow1981}
\CSLLeftMargin{52. }%
\CSLRightInline{Rohsenow DJ, Marlatt GA. The balanced placebo design: Methodological considerations. \emph{Addictive Behaviors}. 1981;6(2):107-122.}

\bibitem[\citeproctext]{ref-iche9r1}
\CSLLeftMargin{53. }%
\CSLRightInline{International Council for Harmonisation of Technical Requirements for Pharmaceuticals for Human Use. {ICH E9(R1)} addendum on estimands and sensitivity analysis in clinical trials to the guideline on statistical principles for clinical trials. Published online 2019.}

\bibitem[\citeproctext]{ref-senn2018}
\CSLLeftMargin{54. }%
\CSLRightInline{Senn S. Statistical pitfalls of personalized medicine. \emph{Nature}. 2018;563(7733):619-621.}

\bibitem[\citeproctext]{ref-dazaCausalAnalysisSelftracked2018}
\CSLLeftMargin{55. }%
\CSLRightInline{Daza EJ. Causal analysis of self-tracked time series data using a counterfactual framework for {N-of-1} trials. \emph{Methods of Information in Medicine}. 2018;57(S 01):e10-e21. doi:\href{https://doi.org/10.3414/ME16-02-0044}{10.3414/ME16-02-0044}}

\bibitem[\citeproctext]{ref-gaertnerCarryover2023}
\CSLLeftMargin{56. }%
\CSLRightInline{Gärtner T, Schneider J, Arnrich B, Konigorski S. Comparison of {Bayesian} networks, {G}-estimation and linear models to estimate causal treatment effects in aggregated {N-of-1} trials with carry-over effects. \emph{BMC Medical Research Methodology}. 2023;23:191. doi:\href{https://doi.org/10.1186/s12874-023-02012-5}{10.1186/s12874-023-02012-5}}

\bibitem[\citeproctext]{ref-morris2019simulation}
\CSLLeftMargin{57. }%
\CSLRightInline{Morris TP, White IR, Crowther MJ. Using simulation studies to evaluate statistical methods. \emph{Statistics in Medicine}. 2019;38(11):2074-2102.}

\end{CSLReferences}

\end{document}
